\setcounter{page}{1}
\chapter*{幂、指、对函数} % {第4章校本作业}

\section*{幂函数\&指数幂的拓展（1）}
\subsection*{指数幂的拓展}

\begin{enumerate}
\item 当$x<0$时, $\abs{x}+\sqrt[6]{x^6}+2\sqrt[3]{x^3}+3\sqrt[4]{x^4}$=\blkda{$-3x$}. 

\item 若$a>0$, $b>0$, 化简: $\left(a^{2-\sqrt{3}}b\right)^{2+\sqrt{3}}\cdot b^{2-\sqrt{3}}$=\blkda{$ab^4$}. 
\begin{jieda}
$E=(a^{2-\sqrt{3}})^{2+\sqrt{3}}\cdot b^{(2+\sqrt{3})+(2-\sqrt{3})}=a^{(2-\sqrt{3})(2+\sqrt{3})}\cdot b^4=ab^4$. 
\end{jieda}

\item 化简:
\begin{enumerate}
\item $\dfrac{5x^{-\frac{2}{3}}y^{\frac{1}{2}}}{\left(\frac{1}{2}x^{-1}y^{\frac{1}{2}}\right)\left(\frac{5}{8}x^{\frac{1}{3}}y^{\frac{1}{6}}\right)}
\times \sqrt[3]{\dfrac{x^2}{\sqrt{y}}}$; 

\item $\left(x^{\frac{\sqrt{3}}{2}}\right)^{\sqrt{3}}\cdot\sqrt{x}(x>0)$; 

\item $\dfrac{\left(a^{\frac{2}{3}}b^{\frac{1}{2}}\right)\left(-3a^{\frac{1}{2}}b^{\frac{1}{3}}\right)}{\frac{1}{3}a^{\frac{1}{6}}b^{\frac{5}{6}}}(a>0,b>0)$.
\end{enumerate}
\begin{jieda}
(1) $\dfrac{5x^{-\frac{2}{3}}y^{\frac{1}{2}}}{\left(\frac{1}{2}x^{-1}y^{\frac{1}{2}}\right)\left(\frac{5}{8}x^{\frac{1}{3}}y^{\frac{1}{6}}\right)}
\times \sqrt[3]{\dfrac{x^2}{\sqrt{y}}} = \dfrac{5x^{-\frac{2}{3}}y^{\frac{1}{2}}}{\frac{5}{16}(x^{-\frac{2}{3}})y^{\frac{2}{3}}} \times \dfrac{x^{\frac{2}{3}}}{y^{\frac{1}{6}}} = 16y^{-\frac{1}{6}} \times y^{-\frac{1}{6}} \times x^{\frac{2}{3}} = 16x^{\frac{2}{3}}y^{-\frac{1}{3}}$

(2) $x^2$; 

(3) $-9a$. 
\end{jieda}


\item 已知 $x+y=12, x y=9$, 且 $x<y$, 求 $\dfrac{x^{\frac{1}{2}}+y^{\frac{1}{2}}}{x^{\frac{1}{2}}-y^{\frac{1}{2}}}$.
\begin{jieda}
    $\dfrac{x^{\frac{1}{2}}+y^{\frac{1}{2}}}{x^{\frac{1}{2}}-y^{\frac{1}{2}}} = \dfrac{\left ( x^{\frac{1}{2}}+y^{\frac{1}{2}} \right) \cdot \left ( x^{\frac{1}{2}}+y^{\frac{1}{2}} \right)}{\left ( x^{\frac{1}{2}}-y^{\frac{1}{2}} \right) \cdot \left ( x^{\frac{1}{2}}+y^{\frac{1}{2}} \right)} = \dfrac{x+y+2\sqrt{xy}}{x-y} = \dfrac{12 + 6}{x-y} = \dfrac{18}{-\sqrt{(x+y)^2-4xy}} = \dfrac{18}{-\sqrt{108}} = -\sqrt{3}$
\end{jieda}
\begin{vspace_}
\vsp[2]
\end{vspace_}

\item 已知 $4^x=a$, 求 $\dfrac{2^{3 x}+2^{-3 x}}{2^x+2^{-x}}$ 的值.
\begin{jieda}
    因为$4^x=a$，所以$2^x = a^{\frac{1}{2}}$，故$\dfrac{2^{3 x}+2^{-3 x}}{2^x+2^{-x}} = \dfrac{a^{\frac{3}{2}} + a^{-\frac{3}{2}}}{a^{\frac{1}{2}} + a^{-\frac{1}{2}}} = a + \dfrac{1}{a} - 1$ \\
    \ps 也可以写做$\dfrac{a^3+1}{a^2+a}$
\end{jieda}
\begin{vspace_}
\vsp[2]
\end{vspace_}

\end{enumerate}

\subsection*{幂函数}

A组：
\begin{enumerate}
\item 若$f(x)=ax^2+b$是幂函数, 则实数$a,b$满足条件\blkda[2]{$a=1,b=0$}. 

\item 若幂函数的图像经过点 $(4,2)$, 则此幂函数为\blkda[2]{$y = x^{\frac{1}{2}}$}.

\item 幂函数$y=x^{-3}$在区间$[-2,-1]$上的值域为\blkda{$[-1,-\frac{1}{8}]$}. 
\begin{jieda}
$y=x^{-3}$在$(-\infty,0)$和$(0,+\infty)$都严格单调递减. 
\end{jieda}

\item 幂函数 $y=x^a$ 的图像不经过原点, 那么 $a$ 的取值范围\blkda[2]{$a \leq 0$}.

\item 如果幂函数 $y=x^a$ 的图像, 当 $0<x<1$ 时, 在直线 $y=x$ 的上方, 那么 $a$ 的取值范围是\blkda[2]{$a < 1$}.

\item 若一个幂函数的图像经过点 $\left(3, \frac{1}{9}\right)$, 则它的单调减区间是\blkda[2]{$(0, +\infty)$}.
\begin{jieda}
    该幂函数为$y = x^{-2}$
\end{jieda}

\item 设 $a \in\left\{-2,-\frac{1}{2},-\frac{1}{3}, \frac{1}{3}, \frac{1}{2}, 2\right\}$, 已知函数 $y=x^a$ 为偶函数, 且在 $(0,+\infty)$ 上递减,则 $a$ 可取的值为\blkda[2]{-2}.

\end{enumerate}

B组：
\begin{enumerate}
\item 若幂函数$y=x^k$的图像经过点$(8,4)$, 则函数$y=x^k$的值域是\blkda{$[0,+\infty)$}. 
\begin{jieda}
$8^k=4\iff 2^{3k}=2^2\iff k=\frac{2}{3}$. \\
于是$y=x^{\frac{2}{3}}=\sqrt[3]{x^2}$的值域为$[0,+\infty)$. 
\end{jieda}

\item 函数$f(x)=\begin{cases}
-x^{-1}, & x<0\\
x^{\frac{1}{2}}, & x>0
\end{cases}$, 
若$f(x)=\frac{1}{3}$, 则$x=$\blkda{$-3$或$\frac{1}{9}$}. 
\begin{jieda}
$-x^{-1}=\frac{1}{3}\land x<0$或$x^{\frac{1}{2}}=\frac{1}{3}\land x>0$, 故$x=-3$或$\frac{1}{9}$. 
\end{jieda}

\item 若幂函数 $f(x)=\left(m^2+m-5\right) x^{m^2-2 m-3}$ 的图象不经过原点, 则 $m$ 的值为 \blkda[2]{$2$}.
\begin{jieda}
    因为$f(x)$是幂函数，所以$m^2+m-5 = 1$，即$m = 2 \mbox{或} -3$ \\
    $m = 2$时，$m^2-2 m-3 = -3$满足题意；$m = -3$时，$m^2-2 m-3 = 12$不满足题意.
\end{jieda}

\item 幂函数$y=x^k$的图像当$0<x<1$时, 在直线$y=x$的上方, 当$x>1$时, 在直线$y=x$的下方, 且$k>0$, 试写出一个符合上述条件的幂函数\blkda{$y=x^{\frac{1}{2}}$}. 

\item 函数$y=x^{\frac{3}{5}}$在区间$[-1,1]$上是\dotfill(\kh{A})
\xx{增函数且是奇函数}{增函数且是偶函数}{减函数且是奇函数}{减函数且是偶函数}


\item 下列四个命题中, 正确的是\blkda{\ding{172}\ding{175}}: 
\begin{dingautolist}{172}
\item $y=x^{-4}$是偶函数, 且在$(0,+\infty)$上是减函数; 
\item $y=x^{\frac{3}{2}}$是奇函数, 且在$(0,+\infty)$上是增函数; 
\item $y=x^{-\frac{1}{2}}$是偶函数, 且在$(0,+\infty)$上是减函数; 
\item $y=x^{-\frac{4}{5}}$是偶函数, 且在$(0,+\infty)$上是减函数. 
\end{dingautolist}
\begin{jieda}
$y=x^{\frac{m}{n}}=\sqrt[n]{x^m}$, 其中$m,n\zz$, 且$(m,n)=1$, $n\neq 0$. 
\end{jieda}

\item 设$y=f(x)$和$y=g(x)$是两个不同的幂函数, 集合$M=\{x\mid f(x)=g(x)\}$, 则集合$M$中的元素个数是\dotfill(\kh{B})
\xx{1或2或0}{1或2或3}{1或2或3或4}{0或1或2或3}
\begin{jieda}
都过$(1,1)$, 故A/D皆不选. \\
对于$x^m=x^n(m,n\zq)$, \\
(1) 若$x=0$有意义, 则$x=0$是根; \\
(2) 当$x\neq 0$时, 可得$x^{\frac{m}{n}}=1$, 故$x$只有$\pm1$的可能. \\
又$y=x$与$y=x^3$恰有三个交点, 故选B.  
\end{jieda}

\item 下列命题中, 真命题的是\dotfill(\kh{C})
\xx{幂函数中不存在既不是奇函数, 又不是偶函数的函数;}
{如果一个幂函数不是偶函数, 那么它一定是奇函数;}
{图像不经过$(-1,1)$的幂函数, 一定不是偶函数;}
{如果两个幂函数有三个公共点, 那么这两个函数一定相同.}


\item 若幂函数$y=x^{m^2-2m-3}(m\zz)$的图像与坐标轴没有交点, 且图像关于$y$轴对称, 求$m$的值. 
\begin{jieda}
$m^2-2m-3 \leq 0\iff m\in[-1,3]$, 又$m\zz$, 故$m=-1,0,1,2,3$. \\
当$m=0$或2时, $m^2-2m-3=-3$, 不符合条件; \\
当$m=1$时, $m^2-2m-3=-4$, 符号条件; \\
当$m=-1$或3时, $m^2-2m-3=0$, 符号条件; \\
综上, $m=-1,1,3$. 
\end{jieda}


\item 已知幂函数$f(x)=(m^2-m-1)x^{m^2+7m+11}$, 问: 当实数$m$取什么值时, 
\begin{enumerate}
\item 函数$f(x)$是正比例函数; 
\item 函数$f(x)$是反比例函数; 
\item 函数$f(x)$是幂函数. 
\end{enumerate}
\begin{jieda}
\begin{enumerate}[(1)]
\item $m^2+7m+11=1\iff m=-5$或$-2$, 此时$m^2-m-1\neq 0$. 故$m=-5$或$-2$. 

\item $m^2+7m+11=-1\iff m=-4$或$-3$, 此时$m^2-m-1\neq 0$. 故$m=-4$或$-3$. 

\item $m^2-m-1=1\iff m=-1$或2. 
\end{enumerate}
\end{jieda}
\begin{vspace_}
\vsp[3]
\end{vspace_}


\item 已知$(a+1)^{-\frac{3}{5}}<(3-2a)^{-\frac{3}{5}}$, 求实数$a$的取值范围. 
\begin{jieda}
\begin{enumerate}
    \item 方法1 \\
    构建函数$f(x) = x^{-\frac{3}{5}}$，原不等式转化为$f(a+1) < f(3-2a)$ \\
    易知函数$f(x) = x^{-\frac{3}{5}}$有两个单调减区间，奇函数，在$x > 0$时$f(x) > 0$，在$x < 0$时$f(x) < 0$ 
    \begin{enumerate}
        \item 若$(a+1)(3-2a) > 0$，即$a \in (-1, \dfrac{3}{2})$，则$a+1$和$3-2a$位于同一个单调区间，此时根据单调性，$f(a+1) < f(3-2a)$被转化为$a+1 > 3-2a$。故$a \in (\dfrac{2}{3}, \dfrac{3}{2})$
        \item 若$a+1< 0 < 3-2a$，即$a \in (-\infty, -1)$，则$a+1$和$3-2a$位于不同单调区间，并且函数的性质可知满足$f(a+1) < f(3-2a)$
    \end{enumerate}
    综上，$a \in (-\infty, -1) \cup (\dfrac{2}{3}, \dfrac{3}{2})$
    \item 方法2 \\
    构建函数$f(x) = x^{\frac{3}{5}}$，原不等式转化为$f(\dfrac{1}{a+1}) < f(\dfrac{1}{3-2a})$ \\
    易知函数$f(x) = x^{\frac{3}{5}}$为$R$上单调递增的奇函数，因此$f(\dfrac{1}{a+1}) < f(\dfrac{1}{3-2a})$即$\dfrac{1}{a+1} < \dfrac{1}{3-2a}$ \\
    移项通分有$\dfrac{2-3a}{(a+1)(3-2a)} < 0$，解集为$a \in (-\infty, -1) \cup (\dfrac{2}{3}, \dfrac{3}{2})$
\end{enumerate}
\end{jieda}
\begin{vspace_}
\vsp[3]
\end{vspace_}


\item 画出下列函数的图象：
\begin{enumerate}
\item $y=x^{\frac{1}{4}}$
\item $y=x^{\frac{3}{5}}$
\item $y=x^{\frac{4}{3}}$
\item $y=x^{-\frac{1}{3}}$
\item $y=x^{-\frac{2}{5}}$
\item $y=x^{-\frac{3}{2}}$
\end{enumerate}
\begin{jieda}
    \begin{center}
    \includegraphics[width=0.3\textwidth]{bodyxb//src/bx01ch04sec1mihanshuB12.png}
\end{center}
\end{jieda}
\begin{vspace_}
\vsp[5]
\end{vspace_}

\item 讨论函数 $y=\left(k^2+k\right) x^{k^2-2 k-1}$ 在 $x>0$ 时随着 $x$ 的增大其函数值的变化情况.

\begin{jieda}
    当$k^2+k$与$k^2-2 k-1$同号的时候，函数严格单调递增，一者为零的时候函数值恒定，异号的时候函数严格单调递减，即
    \begin{enumerate}
        \item $k \in \{0, -1, 1-\sqrt{2}, 1+\sqrt{2}\}$时函数值恒定
        \item $k \in (-\infty, -1) \cup (1-\sqrt{2}, 0) \cup (1+\sqrt{2}, +\infty)$函数值随自变量增大而增大
        \item $k \in (1, 1-\sqrt{2}) \cup (0, 1+\sqrt{2})$函数值随自变量增大而减小
    \end{enumerate}
\end{jieda}

\begin{vspace_}
\vsp[3]
\end{vspace_}


\item 已知幂函数$f(x)=x^{-\frac{1}{2}p^2+p+\frac{3}{2}}(p\zz)$在$(0,+\infty)$上是严格增函数, 且在定义域上是偶函数. 
\begin{enumerate}
\item 求$p$的值, 并写出函数$f(x)$的解析式; 
\item 对于(1)中求得的函数$f(x)$, 设函数$g(x)=-q\cdot f(f(x))+(2q-1)\cdot f(x)+1$, 问是否存在实数$q(q<0)$, 
使得$g(x)$在区间$(-\infty,-4]$上是严格减函数, 且在$(-4,0)$上是严格增函数? 若存在, 请求出$q$; 若不存在, 请说明理由.
\end{enumerate}
\begin{jieda}
\begin{enumerate}[(1)]
\item $-\frac{1}{2}p^2+p+\frac{3}{2}>0\iff p\in(-1,3)$, 又$p\zz$, 故$p=0,1,2$. \\
当$p=0$或2时, $-\frac{1}{2}p^2+p+\frac{3}{2}=\frac{3}{2}$, 不符合条件; \\
当$p=1$时, $-\frac{1}{2}p^2+p+\frac{3}{2}=2$, 符号条件. \\
综上, $p=1$, 此时$f(x)=x^2$. 

\item $g(x)=-qx^4+(2q-1)x^2+1$, 设$g(x) = h(f(x))$，$h(u) = -qu^2 + (2q-1)u + 1, u = f(x) = x^2$ \\
内层函数$f(x)$在$(-\infty, 0)$严格减，在$[0, +\infty)$严格增，值域为$[0， +\infty)$ \\
外层函数$h(u)$在$(-\infty, \dfrac{2q-1}{2q})$严格减，在$[\dfrac{2q-1}{2q}, +\infty)$严格增 \\
由$q < 0$知$\dfrac{2q-1}{2q} > 0$，故当$x \in (-\infty, -\sqrt{\dfrac{2q-1}{2q}}), f(x)$严格减，$u \geq \dfrac{2q-1}{2q}, h(u)$严格增，故$g(x)$严格减; 当$x \in (-\sqrt{\dfrac{2q-1}{2q}}, 0), f(x)$严格减，$0 < u < \dfrac{2q-1}{2q}, h(u)$严格减，故$g(x)$严格增.\\
因此，当$-\sqrt{\dfrac{2q-1}{2q}} = -4$，即$q=-\dfrac{1}{30}$时符合题目要求. 
\end{enumerate}
\end{jieda}
\begin{vspace_}
\vsp[3]
\end{vspace_}
\end{enumerate}

C组：
\begin{enumerate}
    \item 已知函数$f(x)$满足$f(x+1) = 1+\sqrt{2f(x)-f^2(x)}(x \in \mathbf{R})$，求$f(1) + f(2020)$的最大值.
    \begin{jieda}
        易知应分析函数的周期性 \\
        首先分析$f(x)$取值范围，因为$f(x+1) = 1+\sqrt{2f(x)-f^2(x)}$,所以有$2f(x)-f^2(x) \geq 0$，故$f(x) \in [0, 2]$,$f(x+1) \in [1, 2]$，因此$f(x) \in [1, 2]$ \\
        移项后两侧平方有$f^2(x+1) - 2f(x+1) + 1 = 2f(x) - f^2(x)$，若设$g(x) = 2f(x) - f^2(x)$，则$g(x) + g(x+1) = 1$。(利用结论，可以判断$g(x)$周期为2) \\
        进一步迭代有$g(x+1) + g(x+2) = 1$，两式相减有$g(x)= g(x+2)$，故$g(x)$周期为2，即$2f(x)-f^2(x) = 2f(x+2)  - f^2(x+2)$ \\
        移项有$2(f(x)-f(x+2))= (f(x) + f(x+2))(f(x)- f(x+2))$ \\
        由$f(x) \in [1, 2]$知$f(x) + f(x+2) \geq 2$，故$f(x) = 1$或$f(x) = f(x+2)$ \\
        故$f(1) + f(2020) = f(1) + f(0) = f(0) + 1 + \sqrt{2f(0)-f^2(0)}$ \\
        设$t = f(0) (t \in [1, 2]), f(1) + f(2020) = h(t) = t + 1 + \sqrt{2t-t^2} = 2 + (t-1) + \sqrt{1-(t-1)^2}$ 
        \begin{enumerate}
            \item 方法1 \\
            由柯西不等式知$2 + (t-1) + \sqrt{1-(t-1)^2}\leq 2 + 2\sqrt{(t-1)^2 + 2t-t^2} = 2+\sqrt{2}$，当$t = 1+\dfrac{\sqrt{2}}{2}$时取等.
            \item 方法2 \\
            设$k = (t-1)^2 (k \in [0, 1])$，$2 + (t-1) + \sqrt{1-(t-1)^2} = 2 + \sqrt{k} + \sqrt{1-k}$，设$y = \sqrt{k} + \sqrt{1-k}$，故$y^2 = 1 + 2\sqrt{k(1-k)} \leq 2$，故$2 + \sqrt{k} + \sqrt{1-k} \leq 2 + \sqrt{2}$，当$k = \dfrac{1}{2}$时取等.
        \end{enumerate}
        综上，$f(1) + f(2020) \leq 2+2\sqrt{2}$，当$f(0) = 1+\dfrac{\sqrt{2}}{2}$时取等.
         
    \end{jieda}
    \begin{vspace_}
    \vsp[3]
    \end{vspace_}
\end{enumerate}



\section*{幂函数\&指数幂的拓展（2）}

A组：
\begin{enumerate}
\item 将$\left(\sqrt{x^{\frac{1}{3}}\sqrt[3]{x^{-2}}}\right)^{-\frac{8}{5}}$化成分数指数幂为\blkda{$x^{\frac{4}{15}}$}. 

\item 化简$\left(\sqrt[3]{\sqrt[6]{a^9}}\right)^4\cdot\left(\sqrt[6]{\sqrt[3]{a^9}}\right)^4$的结果为\blkda{$a^4$}. 
\begin{jieda}
$a$的指数: $9\cdot\frac{1}{6}\cdot\frac{1}{3}\cdot 4+9\cdot\frac{1}{3}\cdot\frac{1}{6}\cdot4=4$
\end{jieda}

\item $\dfrac{a^{\frac{4}{3}}-8a^{\frac{1}{3}}b}{a^{\frac{2}{3}}+2\sqrt[3]{ab}+4b^{\frac{2}{3}}}\div\left(1-2\sqrt[3]{\dfrac{b}{a}}\right)\times\sqrt[3]{a}$. 
\begin{jieda}
$a$. 
\end{jieda}

\item 已知幂函数 $y=f(x)$ 的图像经过点 $(3, \sqrt{3})$, 则 $f(4)=$ \blkda{$2$}.
\item 已知幂函数 $y=\left(m^2-m-1\right) x^{\frac{m}{3}}$ 的图像关于 $y$ 轴对称, 则 $m=$ \blkda{$2$}.

\end{enumerate}

B组：

\begin{enumerate}
\item 若幂函数$y=x^{\frac{m}{3}}(m\zn)$是奇函数, 则$m$的最小值为\blkda{1}.

\item 若幂函数$y=(m^2-m-1)x^{-5m-3}$在$(0,+\infty)$上为减函数, 则$m$的值为\blkda{2}. 
\begin{jieda}
$m^2-m-1=1\iff m=-1,2$. 分别代入检验即得. 
\end{jieda}

\item 若幂函数$y=x^{\frac{q}{p}}$($\abs{p},\abs{q}$是互质的自然数)的图像关于$y$轴对称, 则$p,q$满足的条件是\blkda[4]{$q$为偶数, $p$为奇数}. 



\item 有关幂函数的下列叙述中, 正确的序号是\blkda{\ding{176}}: 
\begin{dingautolist}{172}
\item 幂函数的图像可能经过第四象限;
\item 两个幂函数的图像至多有两个交点;
\item 若幂函数有增区间, 则该幂函数的指数必是正数;
\item 两个不同的幂函数$f(x)$、$g(x)$, 则函数$h(x)=f(x)\cdot g(x)$还是幂函数;
\item 若幂函数$f(x)$、$g(x)$满足$f^2(x)=g^2(x)$, 则$f(x)=g(x)$. 
\end{dingautolist}
\begin{jieda}
\begin{dingautolist}{172}
\item 当$x>0$时, $y=x^k>0$. 故假. 

\item $y=x$与$y=x^3$有三个交点. 假. 

\item $y=x^{-2}$在$(-\infty,0)$上单调递增. 假.

\item 幂函数的定义域是由``表达式''决定的, 那么$x^3\cdot\dfrac{1}{x}=x^2$. \\
对于$y=x^3\cdot\dfrac{1}{x}$而言, $x\neq 0$; 对于$y=x^2$而言, $x\zr$. \\
那么, 从定义的角度来说, 是假命题. 实际上, 个别点去掉并不影响大局, 故看成真命题未尝不可. 

\item 设$f(x)=x^p$, $g(x)=x^q$, 则$f^2(x)=x^{2p}$, $g^2(x)=x^{2q}$. \\
对于$x=0$而言, $f(x),g(x),f^2(x),g^2(x)$保持一致(要么都为零, 要么都不存在). \\
当$x>0$时, 可得$x^{2p-2q}=1$. 故$2p-2q=0$, 即$p=q$. \\
此时, 若$x<0$有意义, 也是满足条件的. 故真. 
\end{dingautolist}
\end{jieda}



\item 如图曲线是幂函数$y=x^n$在第一象限内的图像, 已知$n$取$\pm2,\pm\frac{1}{2}$四个值, 则曲线$c_1,c_2,c_3,c_4$对应的值依次为\dotfill(\kh{B})
\xx{$-2,-\frac{1}{2},\frac{1}{2},2$}
{$2,\frac{1}{2},-\frac{1}{2},-2$}
{$-\frac{1}{2},-2,2,\frac{1}{2}$}
{$2,\frac{1}{2},-2,-\frac{1}{2}$}
\begin{minipage}{\linewidth}
\centering
\includegraphics[totalheight=3.5cm]{./bodyxb/src/main42chap21mihanshu01.png}
\figcaption*{第5题图}
\end{minipage}


\item 如图是函数$y=x^{\frac{m}{n}}$($m,n\zn$且互质)的图像, 则\dotfill(\kh{C})
\xx{$m,n$是奇数, 且$\dfrac{m}{n}<1$}
{$m$是偶数, $n$是奇数, 且$\dfrac{m}{n}>1$}
{$m$是偶数, $n$是奇数, 且$\dfrac{m}{n}<1$}
{$m$是奇数, $n$是偶数, 且$\dfrac{m}{n}>1$}
\begin{minipage}{\linewidth}
\centering
\includegraphics[totalheight=3.5cm]{./bodyxb/src/main42chap21mihanshu02.png}
\figcaption*{第6题图}
\end{minipage}


\item 函数$g(x)=-\dfrac{x}{x+1}$的图像可以由函数$f(x)=\dfrac{x}{x-1}$的图像经过下列变换得到\dotfill(\kh{B})
\xx{向左平移2个单位, 向上平移2个单位;}
{向左平移2个单位, 向下平移2个单位;}
{向右平移2个单位, 向上平移2个单位;}
{向右平移2个单位, 向下平移2个单位.}
\begin{jieda}
$g(x)=-1+\dfrac{1}{x+1}$, $f(x)=1+\dfrac{1}{x-1}$, 故$g(x)=f(x+2)-2$. 
\end{jieda}


\item 记$\prod\limits_{k=1}^na_k=a_1\cdot a_2\cdot a_3\cdot\cdots\cdot a_n$, 
设$f_n(x)=x^{\frac{n}{3}}, g_n(x)=x^{\frac{n}{5}}, h_n(x)=x^{\frac{n}{2}}$, 
下列论断中, 正确论断的序号为\blkda{\ding{172}}.
\begin{dingautolist}{172}
\item 使函数$F(x)=\prod\limits_{k=1}^nf_k(x)$是偶函数的最小正整数$n$的值等于3;
\item 使函数$G(x)=\prod\limits_{k=1}^ng_k(x)$是偶函数的最小正整数$n$的值等于4;
\item 使函数$H(x)=\prod\limits_{k=1}^nh_k(x)$是奇函数的最小正整数$n$的值等于3.
\end{dingautolist}
\begin{jieda}
积函数的奇偶性. 
\begin{dingautolist}{172}
\item $f_n(x)$为奇函数$\iff n$为奇数, $f_n(x)$为偶函数$\iff n$为偶数, 故真. 
\item $g_n(x)$为奇函数$\iff n$为奇数, $g_n(x)$为偶函数$\iff n$为偶数, 故假. $n_{\min}=3$. 
\item $h_1(x)$的定义域为$[0,+\infty)$, 故非奇非偶. 故假. 
\end{dingautolist}
\end{jieda}


\item 已知幂函数 $f(x)=x^{m^2-2 m-3}(m \in Z)$ 为偶函数, 且在 $(0,+\infty)$ 上是严格减函数.
\begin{enumerate}
\item 求函数 $f(x)$ 的解析式;
\item 设 $F(x)=\sqrt[4]{f(x)}+\frac{1}{\sqrt[4]{f(x)}}$, 试判断函数 $F(x)$ 的奇偶性及单调性.
\end{enumerate}
\begin{jieda}
    \begin{enumerate}
        \item 由$f(x)$在 $(0,+\infty)$ 上是严格减函数知$m^2-2 m-3 < 0$，由$f(x)$为偶函数知$m^2-2 m-3$为偶数. 故$m = 1$，$f(x) = x^{-4}$
        \item $F(x) = \dfrac{1}{|x|} + |x|$，为偶函数，在$(-1, 0), (1, +\infty)$单调增，在$(-\infty, -1), (0, 1)$单调减.
    \end{enumerate}
\end{jieda}
\begin{vspace_}
\vsp[3]
\end{vspace_}


\item 已知$f(x)=\dfrac{x^{\frac{1}{3}}-x^{-\frac{1}{3}}}{5}$, $g(x)=\dfrac{x^{\frac{1}{3}}+x^{-\frac{1}{3}}}{5}$. 
\begin{enumerate}[(1)]
\item 证明$f(x)$是奇函数, 并求$f(x)$的单调区间;
\item 分别计算$f(4)-5f(2)g(2)$和$f(9)-5f(3)g(3)$的值, 由此概括出涉及$f(x)$和$g(x)$的对所有不等于零的实数$x$都成立的一个等式, 并加以证明. 
\end{enumerate}
\begin{jieda}
\begin{enumerate}[(1)]
\item 奇函数易得, $(-\infty,0)$和$(0,+\infty)$单调递增. (注意定义域)
\item $5f(x)g(x)=f(x^2)$. (原式显然是个平方差公式的形式)
\end{enumerate}
\end{jieda}

\begin{vspace_}
\vsp[3]
\end{vspace_}


\item 研究函数$f(x)=\dfrac{x^2+4x+5}{x^2+4x+4}$的单调性, 比较$f(-\pi)$与$f(-\frac{\sqrt{2}}{2})$的大小. 
\begin{jieda}
$f(x)=1+\dfrac{1}{(x+2)^2}$. \\
于是$f(x)$在$(-\infty,-2)$上严格单调递增, 在$(-2,+\infty)$上严格单调递减, \\
且$f(x)$的图像关于$x=-2$对称. \\
故$f(x_1)<f(x_2)\iff\abs{-2-x_1}>\abs{-2-x_2}>0$, 故$f(-\pi)>f(-\frac{\sqrt{2}}{2})$. 
\end{jieda}

\begin{vspace_}
\vsp[3]
\end{vspace_}


\item 已知幂函数$f(x)=x^{-m^2+2m+3}(m\zz)$为偶函数, 且在区间$(0,+\infty)$上是严格增函数. 
\begin{enumerate}[(1)]
\item 求函数$f(x)$的解析式; 
\item 设函数$g(x)=2\sqrt{f(x)}-8x+q-1$, 若不等式$g(x)>0$对任意$x\in[-1,1]$恒成立, 求实数$q$的取值范围. 
\end{enumerate}
\begin{jieda}
\begin{enumerate}[(1)]
\item $-m^2+2m+3>0\iff m\in(-1,3)$. 又$m\zz$, 故$m=0,1,2$. \\
对应的$-m^2+2m+3$依次为$3,4,3$, 故$m=1$, $f(x)=x^4$. 

\item $g(x)=2x^2-8x+q-1>0\iff q>-2x^2+8x+1=-2(x-2)^2+9\in[-9,7]$. \\
故$q>7$. 
\end{enumerate}
\end{jieda}

\begin{vspace_}
\vsp[3]
\end{vspace_}


\item 已知函数$f(x)=x^3+x$. 
\begin{enumerate}[(1)]
\item 指出函数$f(x)$在定义域$\bR$上的奇偶性与单调性(无需证明); 
\item 若$a,b,c\zr$, 且$a+b>0$, $b+c>0$, $c+a>0$, 求证: $f(a)+f(b)+f(c)>0$; 
\item 是否存在自然数$n$, 使$f(n)=1000$?  若存在, 求出$n$; 若不存在, 请说明理由. 
\end{enumerate}
\begin{jieda}
\begin{enumerate}[(1)]
\item 奇函数, 严格单调递增. 
\item $x_1+x_2>0\iff f(x_1)+f(x_2)>0$. 然后$\sum$即可. 
\item $f(9)<1000<f(10)$, 故不存在. 
\end{enumerate}
\end{jieda}

\begin{vspace_}
\vsp[3]
\end{vspace_}

\end{enumerate}

C组：
\begin{enumerate}
    \item 已知$a$是$128$的七次方根，求$\dfrac{1}{1+\sqrt[4]{a}} + \dfrac{1}{1-\sqrt[4]{a}} + \dfrac{2}{1+\sqrt{a}} + \dfrac{4}{1+a}$的值.
    \begin{jieda}
        $\dfrac{1}{1+\sqrt[4]{a}} + \dfrac{1}{1-\sqrt[4]{a}} + \dfrac{2}{1+\sqrt{a}} + \dfrac{4}{1+a} \\
        = \dfrac{(1-\sqrt[4]{a})+(1+\sqrt[4]{a})}{(1+\sqrt[4]{a})(1-\sqrt[4]{a})} + \dfrac{2}{1+\sqrt{a}} + \dfrac{4}{1+a} \\
        = \dfrac{2}{1-\sqrt{a}} + \dfrac{2}{1+\sqrt{a}} + \dfrac{4}{1+a} \\
        =  \dfrac{2(1+\sqrt{a})+2(1-\sqrt{a})}{(1-\sqrt{a})(1+\sqrt{a})} + \dfrac{4}{1+a} \\
        = \dfrac{4}{1-a} + \dfrac{4}{1+a} \\
        = \dfrac{4(1+a) + 4(1-a)}{(1-a)(1+a)} \\
        = \dfrac{8}{1-a^2} \\
        = -\dfrac{8}{3}$
    \end{jieda}
    \begin{vspace_}
    \vsp[3]
    \end{vspace_}
    \item 已知$m$是正整数，幂函数$f(x) = x^{1+m-m^2}$在$(0, +\infty)$上时严格增函数.
    \begin{enumerate}
        \item 求幂函数$f(x)$的解析式
        \item 设$\varphi(x) = \dfrac{a[f(x)]^2+1}{bf(x)+c}$（$a, b, c$为整数），若$\varphi(x)$是奇函数，又$\varphi(1) = 2, \varphi(2) < 3$，求$a, b, c$的值.
    \end{enumerate}
    \begin{jieda}
        \begin{enumerate}
            \item 根据$f(x)$在$(0, +\infty)$上时严格增函数知$1+m-m^2>0$，又因为$m$是正整数，所以$m = 1$,故$f(x) = x$
            \item $\varphi(x) = \dfrac{ax^2+1}{bx+c}$. 因为$\varphi(x)$是奇函数，故$0$不在定义域内，否则$\dfrac{1}{c}$必不为0，因此$c = 0$，经检验，此时$\varphi(x)$为奇函数. \\
            $\varphi(1) = \dfrac{a+1}{b} = 2$，即$b = \dfrac{a+1}{2}$ \\
            $\varphi(2) = \dfrac{4a+1}{2b} = \dfrac{4a+1}{a+1} = 4 - \dfrac{3}{a+1} < 3$，故$a \in \{0, 1\}$ \\
            又因为$b$是整数，故$a = 1$\\
            综上，$a = 1, b = 1, c = 0$
        \end{enumerate}
    \end{jieda}
    \begin{vspace_}
    \vsp[3]
    \end{vspace_}
\end{enumerate}


\section*{指数函数（1）}

A组：
\begin{enumerate}

\item 若函数$f(x)=m^x$在$\bR$上是严格减函数, 则$m$的取值范围是\blkda{$(0,1)$}. 

\item 已知函数$f(x)=a^{x}$的图像恒过定点$P$, 则$P$点坐标是\blkda{$(0,1)$}. 
\begin{jieda}
$a^0=1$!!!
\end{jieda}

\item 如图, 指数函数$y=a^x$, $y=b^x$, $y=c^x$, $y=d^x$, 则$a,b,c,d$与1的大小关系为\dotfill(\kh{B})
\xx{$a<b<1<c<d$}
{$b<a<1<d<c$}
{$1<a<b<c<d$}
{$a<b<1<d<c$}
\begin{minipage}{\linewidth}
\centering
\includegraphics[totalheight=4cm]{./bodyxb/src/main42chap21zhishuhanshu01.png}
\end{minipage}

\item 函数 $f(x)=\sqrt{2^x-1}+(x-2)^0$ 的定义域为\blkda[3]{$[0,2)\cup(2,+\infty)$}.

\item 函数 $y=(2 a-3)^x$ 是指数函数, 则 $a$ 的取值范围是\blkda[3]{$(\frac{3}{2},2)\cup (2,+\infty)$}.

\item 已知函数 $f(x)=(a-2)^x$, 当 $x<0$ 时, $f(x)>1$ 恒成立, 则实数 $a$ 的取值范围是\blkda{$(2, 3)$}.

\item 函数 $y=a^{x-2}+1 \quad(a>0$, 且 $a \neq 1)$ 的图像必经过定点\blkda{$(2, 2)$}.

\item 函数 $f(x)=\left(\frac{1}{3}\right)^x-1, x \in[-1,2]$ 的值域为\blkda{$[-\frac{8}{9}, 2]$} .

\item 设函数 $f(x)=\left\{\begin{array}{ll}2^{-x}, & x \leq 0 \\ 1, & x>0\end{array}\right.$ ，则满足 $f(2 x+1)<f(3 x)$ 的 $x$ 解集为\blkda{$(-\infty, 0)$}.
\begin{jieda}
    原不等式转化为$2x+1 > 3x$ 且 $0 > 3x$，解得$x \in (-\infty, 0)$
\end{jieda}

\end{enumerate}

B组：

\begin{enumerate}

\item 已知函数\ding{172} $y=(-3)^x$, \ding{173} $y=(\frac{1}{3})^x$, \ding{174} $y=3^x$, \ding{175} $y=x^{-3}$, \ding{176} $y=(\sqrt{2})^x$, 
其中是指数函数的有\blkda{\ding{173}\ding{174}\ding{176}}. 

\item 若函数$f(x)=(m^2-1)^x$在$\bR$上是严格减函数, 则$m$的取值范围是\blkda[3.5]{$(-\sqrt{2},-1)\cup(1,\sqrt{2})$}. 
\begin{jieda}
$0<m^2-1<1\iff m\in(-\sqrt{2},-1)\cup(1,\sqrt{2})$. 
\end{jieda}

\item 满足$(a^2+a+2)^x>(a^2+a+2)^{1-x}$(其中$a\zr$)的实数$x$的取值范围是\blkda{$x>\frac{1}{2}$}. 
\begin{jieda}
$a^2+a+2=(a+\frac{1}{2})^2+\frac{7}{4}>1$, 故$x>1-x$. 
\end{jieda}


\item 已知函数$f(x)=4+a^{x-1}$的图像恒过定点$P$, 则$P$点坐标是\blkda{$(1,5)$}. 
\begin{jieda}
$a^0=1$!!!
\end{jieda}


\item 设集合 $P=\{(x, y) \mid y=k, x \in \mathbf{R}\}, Q=\left\{(x, y) \mid y=a^x+1, x \in \mathbf{R}, a>0\right.$ 且 $\left.a \neq 1\right\}$, 已知 $P \cap Q=\varnothing$, 那么实数 $k$ 的取值范围是\blkda{$(-\infty,1]$}.


\item 若函数$y=a^x$($a>0$且$a\neq1$)在区间$[1,2]$上的最大值比最小值大$\dfrac{a}{2}$, 则实数$a$的值为\blkda{$\frac{1}{2},\frac{3}{2}$}. 
\begin{jieda}
因为单调, 故$\abs{a-a^2}=\dfrac{a}{2}$, 解得$a=\frac{1}{2},\frac{3}{2}$. 经检验, 皆符合条件. 
\end{jieda}

\item 若函数$y=\dfrac{a^x}{2a^2-5a+2}$在$\bR$上为严格减函数, 则$a$的取值范围是\blkda[3]{$(0,\frac{1}{2})\cup(1,2)$}. 
\begin{jieda}
$\begin{cases}
a\in(0,1)\\
2a^2-5a+2>0
\end{cases}$
或$\begin{cases}
a>1\\
2a^2-5a+2<0
\end{cases}$, 
解得$a\in(0,\frac{1}{2})\cup(1,2)$. 
\end{jieda}


\item 若函数$f(x)=\dfrac{1}{2^x+1}$, 则该函数在$\bR$上是\dotfill(\kh{A})
\xx{严格单调递减无最小值}
{严格单调递减有最小值}
{严格单调递增无最大值}
{严格单调递增有最大值}
\begin{jieda}
$2^x+1\in(1,+\infty)$, 故$\dfrac{1}{2^x+1}\in(0,1)$. 
\end{jieda}


\item 若$0<a<1$, $b<-1$, 则函数$f(x)=a^x+b$的图像不经过\dotfill(\kh{A})
\xx{第一象限}{第二象限}{第三象限}{第四象限}
\begin{jieda}
$f(0)=1+b<0$(连续性可知必过III,IV). 又$a\in(0,1)$, 故不过I. 
\end{jieda}


\item 函数$y=2^{1-x}+2$的图像可以由函数$y=(\frac{1}{2})^x$的图像经过下面变换得到\dotfill(\kh{C})
\xx{先向左平移1个单位, 再向上平移2个单位}
{先向左平移1个单位, 再向下平移2个单位}
{先向右平移1个单位, 再向上平移2个单位}
{先向右平移1个单位, 再向下平移2个单位}
\begin{jieda}
令$f(x)=(\frac{1}{2})^x=2^{-x}$, 则$y=f(x-1)+2=2^{1-x}+2$. \\
另外, 用$f(0)=1$也是可以的, 关键在于``若你明白了什么是代表元素那就自由的得''. 
\end{jieda}


\item 若函数 $f(x)=\left\{\begin{array}{l}a^x, x \geq 1 \\ \left(2-\frac{a}{2}\right) x+2, x<1\end{array}\right.$ 是 $\mathbf{R}$ 上的严格增函数, 则实数 $a$ 的取值范围是\blkda{$[\frac{8}{3}, 4)$}.
\begin{jieda}
    $f(x)$ 是 $\mathbf{R}$ 上的严格增函数 $\Longleftrightarrow \left \{ \begin{array}{l}
        a > 1  \\
        2 - \dfrac{a}{2} > 0 \\
        a \geq  (2-\dfrac{a}{2}) + 2
    \end{array}\right.$，解得$a \in [\dfrac{8}{3}, 4)$
\end{jieda}


\item 判断下列函数的奇偶性, 并证明. \\
\begin{enumerate*}
\item $y=\dfrac{a^x-a^{-x}}{2},(a>0, a \neq 1)$;
\item $y=x \cdot \dfrac{a^x-1}{a^x+1},(a>0, a \neq 1)$.
\end{enumerate*}
\begin{jieda}
\begin{enumerate*}[(1)]
\item 奇;
\item 偶. 
\end{enumerate*}
\end{jieda}

\item 解不等式: $0.3^{x^2+x+1}>0.3^{-2x^2+5x}$. 
\begin{jieda}
因为$y=0.3^x$在$\bR$上严格单调递减, 故原不等式
$\iff x^2+x+1<-2x^2+5x\iff x\in(\frac{1}{3},1)$. 
\end{jieda}


\item 解关于$x$的不等式: $a^{x^2}>(\frac{1}{a})^{-2x-1}$(其中$a>0$且$a\neq 1$). 
\begin{jieda}
$a>1$时, 等价于$x^2>2x+1$, 即$x\in(-\infty,1-\sqrt{2})\cup(1+\sqrt{2},+\infty)$; \\
$a\in(0,1)$时, 等价于$x^2<2x+1$, 即$x\in(1-\sqrt{2},1+\sqrt{2})$. \\
故$a\in(0,1)$时, 解集为$(1-\sqrt{2},1+\sqrt{2})$; \\
$a>1$时, 解集为$(-\infty,1-\sqrt{2})\cup(1+\sqrt{2},+\infty)$. 
\end{jieda}


\item 已知函数 $f(x)=4^x+a \cdot 4^{-x}$ 是偶函数.
\begin{enumerate}
\item 求 $a$ 的值;
\item 证明: 对任意实数 $x_1$ 和 $x_2$, 都有 $\frac{1}{2}\left[f\left(x_1\right)+f\left(x_2\right)\right] \geq f\left(\frac{x_1+x_2}{2}\right)$.
\end{enumerate}

\begin{jieda}
\begin{enumerate}
\item $a=1$;
\item 任取 $x_1, x_2$,
$$
\begin{aligned}
& \frac{1}{2}\left[f\left(x_1\right)+f\left(x_2\right)\right]=\frac{1}{2}\left(4^{x_1}+4^{-x_1}+4^{x_2}+4^{-x_2}\right) \geq \frac{1}{2}\left(2 \sqrt{4^{x_1} 4^{x_2}}+2 \sqrt{4^{-x_1} 4^{-x_2}}\right) \\
& =4^{\frac{x_1+x_2}{2}}+4^{-\frac{x_1+x_2}{2}}=f\left(\frac{x_1+x_2}{2}\right)
\end{aligned}
$$
\end{enumerate}
\end{jieda}


\item 设函数 $f(x)=\left(a^2+a-5\right) a^x$ 是指数函数, 将 $f(x)$ 的图像向左平移 1 个单位,再向上平移 2 个单位得到 $y=g(x)$ 的图像.
\begin{enumerate}
    \item 写出 $g(x)$ 的解析式;
    \item 对任意的 $x \in[-1,2], g(x) \geq 3+\frac{m}{2}-\frac{m^2}{2}$ 恒成立,求实数 $m$ 的取值范围.
\end{enumerate}
\begin{jieda}
\begin{enumerate}
\item $g(x)=2^{x+1}+2$
\item $\because g(x)$ 在 $[-1,2]$ 上是严格增函数, $\therefore g(x)_{\min }=g(-1)=3$
$\therefore 3+\frac{m}{2}-\frac{m^2}{2} \leq 3, \therefore m \geq 1$ 或 $m \leq 0$.
\end{enumerate}
\end{jieda}

\item 定义在 $D$ 上的函数 $f(x)$, 如果满足: 对任意 $x \in D$, 存在常数 $M>0$, 都有 $-M \leq f(x) \leq M$成立, 则称 $f(x)$ 是 $D$ 上的有界函数, 其中 $M$ 称为函数 $f(x)$ 的上界. 已知 $f(x)=4^x+a \cdot 2^x-2$.
\begin{enumerate}
\item 当 $a=-2$ 时, 求函数 $f(x)$ 在 $(0,+\infty)$ 上的值域, 并判断函数 $f(x)$ 在 $(0,+\infty)$ 上是否为有界函数，请说明理由；
\item 若函数 $f(x)$ 在 $(-\infty, 0)$ 上是以 2 为上界的有界函数, 求实数 $a$ 的取值范围.
\end{enumerate}
\begin{jieda}
\begin{enumerate}
\item $f(x)=4^x-2 \cdot 2^x-2=\left(2^x-1\right)^2-3$, 函数 $f(x)$ 在 $(0,+\infty)$ 上的值域为 $(-3,+\infty)$; 因此不存在常数 $M>0$, 使得 $-M \leq f(x) \leq M$ 成立,所以 $f(x)$ 在 $(0,+\infty)$ 上不是有界函数;
\item  $-2 \leq 4^x+a \cdot 2^x-2 \leq 2$ 对任意的 $x \in(-\infty, 0)$ 恒成立,

即 $-2^x \leq a \leq \frac{4}{2^x}-2^x$ 对任意的 $x \in(-\infty, 0)$ 恒成立,
因为 $x<0$, 所以 $2^x \in(0,1),-2^x \in(-1,0)$;
又 $y=\frac{4}{2^x}-2^x$ 在 $(-\infty, 0)$ 上是严格减函数, 所以 $y=\frac{4}{2^x}-2^x>\frac{4}{2^0}-2^0=3$,
因此为使 $-2^x \leq a \leq \frac{4}{2^x}-2^x$ 对任意的 $x \in(-\infty, 0)$ 恒成立, 只需 $0 \leq a \leq 3$.
\end{enumerate}

\end{jieda}



\end{enumerate}


\section*{指数函数（2）}

A组：

\begin{enumerate}
\item 若$f(\pi^x)=2x-1$, 则$f(1)=$\blkda{$-1$}. 
\begin{jieda}
$f(\pi^0)=2\cdot0-1$, 故$f(1)=-1$. 
\end{jieda}


\item 函数$y=(\frac{1}{3})^{-2x^2-8x+1}$单调递增区间是\blkda{$[-2,+\infty)$}. 


\item 函数$y=x^2$与函数$y=2^x$的图像的交点个数\blkda{3个}. 
\begin{jieda}
$x\in(-\infty,0)$上有一个交点, 这可以用单调性解决. \\
$x>0$时, $x=2,4$是交点. 
\end{jieda}

\item 已知指数函数 $y=(2-a)^x$ 是严格增函数, 则实数 $a$ 的取值范围是\blkda{$(-\infty, 1)$}.


\item 实数 $a$ 满足 $4^{a-1}<\left(\frac{1}{2}\right)^{a-4}$, 则实数 $a$ 的取值范围为\blkda{$(-\infty, 2)$}.
\begin{jieda}
    原不等式式等价于$2^{2a-2}<2^{4-a}$，由指数函数的单调性知原不等式等价于$2a-2 < 4-a$，解得$a \in (-\infty, 2)$
\end{jieda}

\item 若函数 $y=a^x(a>0$, 且 $a \neq 1)$, 在 $[2,3]$ 上的最大值比最小值大 $\frac{a^2}{2}$, 则 $\mathrm{a}=$ \blkda{$\frac{1}{2},\frac{3}{2}$}.
\begin{jieda}
因为单调, 故$\abs{a^2-a^3}=\dfrac{a^2}{2}$, 解得$a=\frac{1}{2},\frac{3}{2}$. 经检验, 皆符合条件. 
\end{jieda}

\end{enumerate}


B组：

\begin{enumerate}
\item 若函数$f(x)=\dfrac{1}{2^x-1}+a$为奇函数, 则实数$a=$\blkda{$\dfrac{1}{2}$}. 
\begin{jieda}
$f(x)=\dfrac{1}{a^x-1}+\lambda$为奇函数$\iff\lambda=\dfrac{1}{2}$, \\
$g(x)=\dfrac{1}{a^x+1}+\lambda$为奇函数$\iff\lambda=-\dfrac{1}{2}$. 
\end{jieda}


\item 函数$f(x)=\dfrac{3^x}{9^x+1}-\dfrac{1}{2}$的值域是\blkda{$\left(-\dfrac{1}{2},0\right]$}. 
\begin{jieda}
因为$y=3^x+3^{-x}\in[2,+\infty)$, \\
故$f(x)=\dfrac{1}{3^x+3^{-x}}-\dfrac{1}{2}\in\left(-\dfrac{1}{2},0\right]$. 
\end{jieda}

\item 函数$y=3^{\sqrt{x^2-5x+4}}$的递减区间是\blkda{$(-\infty,1]$}, 值域是\blkda{$[1,+\infty)$}. 
\begin{jieda}
定义域$x\in(-\infty,1]\cup[4,+\infty)$. \\
根据复合函数的单调区间的原则``同增异减''可得递减区间是$(-\infty,1]$. \\
$\sqrt{x^2-5x+4}\in[0,+\infty)$, 故$y\in[1,+\infty)$. 
\end{jieda}

\item 已知实数满足等式$\left(\dfrac{1}{2}\right)^a=\left(\dfrac{1}{3}\right)^b$, 下列5个关系式中, 不可能成立的有\blkda{\ding{174}\ding{175}}: \\
\ding{172} $0<b<a$; \ding{173} $a<b<0$; \ding{174} $0<a<b$; \ding{175} $b<a<0$; \ding{176} $a=b$. 
\begin{jieda}
作图: $y=(\frac{1}{2})^x, y=(\frac{1}{3})^x$, 作直线$y=t$, 分别取三种情形: $t\in(0,1),1,(1,+\infty)$. 
\end{jieda}

\item 已知$f(x)=\begin{cases}
(2-a)x+1, & x<1\\
a^x, & x\geq 1
\end{cases}(a>0, a\neq1)$是$\bR$上的严格增函数, 则实数$a$的取值范围是\blkda{$[\frac{3}{2},2)$}. 
\begin{jieda}
$2-a >  0\land a>1\land (2-a)\cdot 1+1\leq a^1$, 解得$a\in[\frac{3}{2},2)$. \\
\ps 单调与严格单调, 各段单调, 区间端点. 
\end{jieda}

\item 函数$y=2^{-\abs{x-1}}$单调减区间是\blkda{$[1,+\infty)$}. 

\item 若函数$f(x)=a^x-x-a(a>0,a\neq 1)$有两个零点, 则实数$a$的取值范围是\blkda{$a>1$}. 
\begin{jieda}
$f(x)=0\iff a^x-a=x$. \\
当$a\in(0,1)$时, 左侧函数单调递减, 右侧单调递增, 故至多一解. \\
当$a>1$时, 左右两侧的函数的图像如下: \\
\begin{minipage}{\linewidth}
\centering
\includegraphics[totalheight=3.5cm]{./bodyxb/src/main42chap21zhishuhanshu03.png}
\end{minipage}
\end{jieda}

\item 若$a>0$且$a\neq 1$, 函数$y=\abs{a^x-1}$与$y=2a$的图像有两个交点, 则$a$的取值范围是\blkda{$(0,\frac{1}{2})$}. 
\begin{jieda}
画图(无所谓$a$与1的大小)易得$2a\in(0,1)\iff a\in(0,\frac{1}{2})$. 
\end{jieda}


\item 要使函数$y=1+2^x+a\cdot4^x$在$x\in\left(-\infty,-1\right]$上$y\geq0$恒成立, 则$a$的取值范围是\blkda{$a\geq-6$}. 
\begin{jieda}
$\iff a\geq\dfrac{-1-2^x}{4^x}=-t^2-t$, 其中$t=\frac{1}{2^x}\in[2,+\infty)$. \\
于是$a\geq-6$. 
\end{jieda}


\item 设函数$f(x)$定义在实数集上, 且图像关于直线$x=1$对称, 当$x\geq 1$时, $f(x)=3^x-1$, 则\dotfill(\kh{B})
\xx{$f(\frac{1}{3})<f(\frac{3}{2})<f(\frac{2}{3})$}
{$f(\frac{2}{3})<f(\frac{3}{2})<f(\frac{1}{3})$}
{$f(\frac{2}{3})<f(\frac{1}{3})<f(\frac{3}{2})$}
{$f(\frac{3}{2})<f(\frac{2}{3})<f(\frac{1}{3})$}
\begin{jieda}
$f(\frac{1}{3})=f(\frac{5}{3})$, $f(\frac{2}{3})=f(\frac{4}{3})$. \\
$f(x)$在$[1,+\infty)$上严格单调递增, $\frac{4}{3}<\frac{3}{2}<\frac{5}{3}$, 于是$f(\frac{4}{3})<f(\frac{3}{2})<f(\frac{5}{3})$. 
\end{jieda}

\item 若函数$y=2^{-x+1}+m$的图像不经过第一象限, 则$m$的取值范围是\dotfill(\kh{A})
\xx{$m\leq -2$}
{$m\geq -2$}
{$m\leq -1$}
{$m\geq -1$}
\begin{jieda}
画图, 于是$f(0)\leq 0\iff m\leq -2$. 
\end{jieda}


\item 若函数$f(x)=\dfrac{a-\me^x}{1+a\cdot\me^x}$($a$为常数)在定义域上为奇函数, 则$a$的值为\dotfill(\kh{D})
\xx{$0$}
{$1$}
{$-1$}
{$1$或$-1$}
\begin{jieda}
首选解法：$f(0)=0$或无意义，即$\dfrac{a-1}{1+a} = 0$或$\dfrac{a-1}{1+a}$无意义，故$a = \pm 1$\\
也可以让函数在$\pm \infty$时的极限互为相反数求解：$x \rightarrow +\infty, f(x) \rightarrow -\dfrac{1}{a}, x \rightarrow -\infty, f(x) \rightarrow a$，由$a = \dfrac{1}{a}$得$a = \pm 1$. \\
把握定义求解有：$f(x) = \dfrac{a-\me^x}{1+a\cdot\me^x}, f(-x) = \dfrac{a-\me^{-x}}{1+a\cdot\me^{-x}} = \dfrac{a\me^{x}-1}{a+\me^{x}}$ \\
由$f(x) = -f(-x)$知$\dfrac{a-\me^x}{1+a\cdot\me^x} = \dfrac{1-a\me^{x}}{a+\me^{x}}$，交叉相乘得到$a^2 - e^{2x} = 1-a^2e^{2x}$，移项整理得$(a^2-1) = (1-a^2)\cdot(e^{2x})$，因为此式恒成立，所以必有$a = \pm 1$
\end{jieda}


\item 若$f(x)$在$(0,+\infty)$上是减函数, 而$f(a^x)$在$\bR$上是增函数, 则实数$a$的取值范围是\dotfill(\kh{A})
\xx{$(0,1)$}
{$(0,1)\cup(1,+\infty)$}
{$(0,+\infty)$}
{$(1,+\infty)$}
\begin{jieda}
$a^x\in(0,+\infty)$, 故根据复合函数的单调性可知$0<a<1$. 
\end{jieda}


\item 已知函数$y=4^x-3\cdot2^x+3$的值域为$\left[1,7\right]$, 则$x$的范围可以是\dotfill(\kh{D})
\xx{$\left[2,4\right]$}
{$(-\infty,0)$}
{$(0,1)\cup\left[2,4\right]$}
{$\left(-\infty,0\right]\cup\left[1,2\right]$}
\begin{jieda}
$t^2-3t+3=1\iff t=1$或2, $t^2-3t+3=7\iff t=-1$或4. \\
由于$2^x\in(0,+\infty)$, 故结合$y=t^2-3t+3$的图像及选项可知, 必须要有$2^x\in[2,4]\iff x\in[1,2]$. \\
至于$[-1,1]$, $2^x$不能遍历, 所以$x\to-\infty$方向是随意的. 
\end{jieda}


\item 已知函数$y=\dfrac{1}{4-2^x}$的图像关于点$P$对称, 则点$P$的坐标为\dotfill(\kh{C})
\xx{$(2,\frac{1}{2})$}
{$(2,\frac{1}{4})$}
{$(2,\frac{1}{8})$}
{$(0,0)$}
\begin{jieda}
设$P(a,b)$, 因为$x\neq 2$, 故$a=2$. \\
$2b=f(1)+f(3)$, 于是$b=\frac{1}{8}$. 
\end{jieda}



\item 已知函数$f(x)=2^x-2$, 则函数$y=\abs{f(\abs{x})}$的图像只可能是\dotfill(\kh{A})
\begin{minipage}{\linewidth}
\centering
\includegraphics[totalheight=3cm]{./bodyxb/src/34713785651095890.png}
\end{minipage}

\begin{jieda}

\begin{minipage}{\linewidth}
\centering
\includegraphics[totalheight=3cm]{./bodyxb/src/main42chap21zhishuhanshu04.png}
\end{minipage}
\end{jieda}



\item 求函数$y=3\cdot4^x-20\cdot2^x+3\cdot4^{-x}-20\cdot2^{-x}$的最小值. 
\begin{jieda}
令$t=2^x+2^{-x}\in[2,+\infty)$, 则
\begin{align*}
y
& = 3(4^x+4^{-x})-20(2^x+2^{-x})\\
& = 3(2^x+2^{-x})^2-20(2^x+2^{-x})-6\\
& = 3t^2-20t-6\\
& = 3(t-\frac{10}{3})^2-\frac{118}{3}
\end{align*}
    故$y_{\min}=-\frac{118}{3}$, 当$t=\frac{10}{3}$即$2^x = \frac{1}{3}\mbox{或}3, x=log_23 \mbox{或}-log23$时取到. 
\end{jieda}


\item 设函数$f(x)=\dfrac{a(a^{2x}-1)}{a^x(a^2-1)}$($a>0$且$a\neq 1$), 证明: $y=f(x)$是$\bR$上为严格增函数. 
\begin{jieda}
$f(x)=\dfrac{a}{a^2-1}\cdot(a^x-a^{-x})$. \\
任取实数集内的$x_1, x_2, x_1 < x_2$ \\
$f(x_2)-f(x_1)=\dfrac{a}{a^2-1}\left(a^{x_2}-a^{x_1} - (a^{-x_2} - a^{-x_1})\right)$. 
\begin{enumerate}
    \item 当$0 < a < 1$时,$a^{x_2}-a^{x_1} < 0$，$a^{-x_2} - a^{-x_1} > 0$，$\dfrac{a}{a^2-1} < 0$，故$f(x_2)-f(x_1) > 0$
    \item 当$a > 1$时,$a^{x_2}-a^{x_1} > 0$，$a^{-x_2} - a^{-x_1} < 0$，$\dfrac{a}{a^2-1} > 0$，故$f(x_2)-f(x_1) > 0$
\end{enumerate}
因此$f(x)$是$\bR$上为严格增函数. \\
\ps $f(x)$并非恒正或者恒负，因此不应使用作商法. \\
指数函数的单调性证明如下：
任取实数集内的$x_1, x_2, x_1 < x_2$，由$a^{x_1} > 0, a^{x_2} > 0$
\begin{enumerate}
    \item 若$a > 1$，则$\dfrac{a^{x_2}}{a^{x_1}} = a^{x_2 - x_1}$，由幂的基本不等式，因为$x_2-x_1 > 0$，故$a^{x_2 - x_1} > 1$，因此$a^x$单调递增.
    \item 若$1 > a > 0$，则$\dfrac{a^{x_2}}{a^{x_1}} = a^{x_2 - x_1} = \dfrac{1}{(\frac{1}{a})^{x_2 - x_1}}$，由幂的基本不等式，因为$x_2-x_1 > 0$，故$(\frac{1}{a})^{x_2 - x_1} > 1$，故$\dfrac{a^{x_2}}{a^{x_1}} < 1$，因此$a^x$单调递减.
\end{enumerate}
\end{jieda}

\item 若$a^{2x}+\frac{1}{2}a^x-\frac{1}{2}\leq 0$($a>0$且$a\neq 1)$, 求函数$y=2a^{2x}-3a^x+4$的值域. 
\begin{jieda}
$a^{2x}+\frac{1}{2}a^x-\frac{1}{2}\leq 0\land a^x>0\iff a^x\in(0,\frac{1}{2}]$. \\
$y=2a^{2x}-3a^x+4=2(a^x-\frac{3}{4})^2+\frac{23}{8}\in[3,4)$. 
\end{jieda}


\item 利用函数图像, 讨论方程$2^x-\abs{x+2}=0$解的个数, 并指出每个解所在区间$[n,n+1)$, 其中$n\zz$. 
\begin{jieda}
需要你的计算器, 需要你的细心, 召唤ggb. \\
\begin{minipage}{\linewidth}
\centering
\includegraphics[totalheight=4cm]{./bodyxb/src/main42chap21zhishuhanshu05.png}
\end{minipage}
于是有3个解, $n=-3,-2,2$. 
\end{jieda}


\item 已知函数$f(x)=\dfrac{2^x}{a}-\dfrac{a}{2^x}(a>0)$是奇函数. 
\begin{enumerate}[(1)]
\item 求$a$的值; 
\item 判断函数$f(x)$在$(-\infty,+\infty)$上的单调性并证明. 
\end{enumerate}
\begin{jieda}
\begin{enumerate}[(1)]
\item $f(0)$有意义, 故$f(0)=0$, 解得$a=\pm1$. \\
又$a>0$, 故$a=1$. 检验后发现符合条件. 
\item 当$a=1$时, $f(x)$单调递增. 
\end{enumerate}
\end{jieda}


\item 已知$f(x)=\dfrac{10^x-10^{-x}}{10^x+10^{-x}}$.
\begin{enumerate}[(1)]
\item 判断函数的奇偶性;
\item 证明$f(x)$在区间$(-\infty ,+\infty )$上是增函数;
\item 求函数$f(x)$的值域.
\end{enumerate}
\begin{jieda}
$f(x)=\dfrac{100^x-1}{100^x+1}$. 
\begin{enumerate}[(1)]
\item $x\zr$, 奇函数. 其实$y=\dfrac{a^x-1}{a^x+1}$, $y=\dfrac{a^x+1}{a^x-1}$都是奇函数. 
\item 定义去证, 最后一定是归结为$y=100^x$的单调性. 
\item $f(x)\in(-1,1)$. 
\end{enumerate}
\end{jieda}


\item 已知$x,y\zr$, 且$2^x+3^y>2^{-y}+3^{-x}$, 求证: $x+y>0$. 
\begin{jieda}
令$f(t)=2^t-3^{-t}$, 则$f(t)$在$\bR$上严格单调递增. \\
条件可化为$2^x-3^{-x}>2^{-y}-3^y$, 即$f(x)>f(-y)$, \\
故$x>-y$, 即$x+y>0$. 
\end{jieda}


\item 已知函数$f(x)=\dfrac{a\cdot2^x+a^2-2}{2^x-1}(x\neq0)$, 其中$a$为实数. 
\begin{enumerate}[(1)]
\item 若函数$f(x)$为奇函数, 求出$a$的值;
\item 若函数$f(x)$在$(0,+\infty )$上是增函数, 求$a$的取值范围;
\item 当$x\in\left[1,2\right]$时, $f(x)>1$恒成立, 求$a$的取值范围. 
\end{enumerate}
\begin{jieda}
\begin{enumerate}
\item $f(1)+f(-1)=0$, 可得$a=-1$或2. 再证一下. \\
当$a=-1$时, $f(x)=-\dfrac{2^x+1}{2^x-1}$; 当$a=2$时, $f(x)=2\cdot\dfrac{2^x+1}{2^x-1}$. \\
\ps $f(-\infty)+f(+\infty)=0$, 可得$a=-1$或2. 

\item $f(x)=a+\dfrac{a^2+a-2}{2^x-1}$, 故$a^2+a-2\leq 0$, 即$a\in[-2,1]$. 

\item $x\in\left[1,2\right]$时, $2^x-1\in[1,3]$, 于是$f(x)>1\iff a\cdot2^x+a^2-2>2^x-1$. \\
即$(a-1)\cdot 2^x+(a^2-1)>0$. \\
于是$\begin{cases}
(a-1)\cdot 2^1+(a^2-1)>0\\
(a-1)\cdot 2^2+(a^2-1)>0
\end{cases}$, 解得$a\in(-\infty,-5)\cup(1,+\infty)$. 
\end{enumerate}
\end{jieda}


\end{enumerate}


\section*{反函数（1）}
\zsd 反函数的概念、求反函数

A组：

\begin{enumerate}
\item 函数$y=2x+1(x\geq1)$的反函数是\blkda[3]{$y=\dfrac{x-1}{2}(x\geq 3)$}. 

\item 函数$y=x^{-\frac{1}{2}}$的反函数是\blkda[3]{$y=x^{-2}(x>0)$}. 

\item 函数$f(x)=\dfrac{2}{1-x^2}(x<-1)$, 则$f^{-1}(-\frac{2}{3})=$\blkda{$-2$}.  

\item ``函数$f(x)(x\zr)$存在反函数''是``函数$f(x)$在$\bR$上为严格增函数''的\blkda[2]{必要不充分} 条件. 


\item 已知$y=f(x)$有反函数, 那么方程$f(x)=a(a$为常数)\dotfill(\kh{C})
\xx{无实数根}
{有且仅有一根}
{至多有一个实根}
{至少有一个实根}


\item 函数$f(x)=x^2-2ax-3$在区间$[1,2]$存在反函数的充分必要条件是\blkda[4]{$a\in(-\infty,1]\cup[2,+\infty)$}.


    
\end{enumerate}

B组：

\begin{enumerate}
\item 函数$f(x)=(x-1)^2+1(x<1)$的反函数为\blkda[4]{$y=1-\sqrt{x-1}(x>1)$}. 

\item 函数$f(x)=\begin{cases}
-x^2-2, & x\geq 0\\
2^{-x}, & x<0
\end{cases}$的反函数的定义域为\blkda[4]{$(-\infty,-2]\cup(1,+\infty)$}. 

\item 已知函数$y=f(x)$是奇函数, 当$x\geq 0$时, $f(x)=3^x-1$, 设函数$f(x)$的反函数是$y=g(x)$, 则$g(-8)=$\blkda{$-2$}. 


\item 下列区间中, 使得$f(x)=2\abs{x}$不存在反函数的区间是\dotfill(\kh{B})
\xx{$[2,4]$}
{$[-4,4]$}
{$[0,+\infty)$}
{$(-\infty,0)$}

\item 已知函数$y=x^2-2ax+3(x\in[1,2]\cup[8,9])$有反函数, 求实数$a$的取值范围. 
\begin{jieda}
$a\in(-\infty,1]\cup[2,\frac{9}{2})\cup(\frac{11}{2},8]\cup[9,+\infty)$. 
\end{jieda}


\item 已知 $f(x)=\frac{a-x}{x-a-1}$, 若函数 $y=f(x)$ 的反函数图像的对称中心是 $(-1,3)$, 则实数 $a$ 的值是\blkda{$2$}.
\begin{jieda}
    由反函数的对称中心知$f(x)$的对称中心为$(3, -1)$，因此$a+1 = 3$，故$a = 2$
\end{jieda}


\item 已知 $f(x)=x^2$, 若函数 $y=f(x), x \in D$ 的值域是 $\{1,4,9\}$, 且函数 $y=f(x)$ 存在反函数, 这样的 $f(x)$ 共有\blkda{$8$}个.
\begin{jieda}
    $\pm 1$二选一，$\pm 2$二选一，$\pm 3$二选一，根据乘法原理，总共$2^3$种选法.
\end{jieda}

\item 已知 $f(x)=\sqrt{x-1}$, 函数 $y=f(x)$ 的反函数为 $y=f^{-1}(x)$, 若 $f^{-1}(x)-a=f(x+a)$有实数根, 则实数 $a$ 的取值范围为\blkda{$[\dfrac{3}{4}, +\infty)$}.
\begin{jieda}
    $f^{-1}(x) = x^2 + 1 (x \geq 0)$. \\
    设$g(x) = f(x+a)$，则$g^{-1}(x) = f^{-1}(x)-a$，因为$g(x)$与$g^{-1}(x)$均为增函数，因此其交点必然在$y = x$上。由图象可知，$a$的取值范围形式应为$[m, +\infty)$，$m$为边界情况对应的值.边界情况为$x^2 + 1 - a = x$只有一解，此时$a = \dfrac{3}{4}$，故$a \in [\dfrac{3}{4}, +\infty)$
\end{jieda}


\item 给出下列几个函数:
\ding{172} $y=x^2-1(x<-2)$; 
\ding{173} $y=x^3-1(x\zr)$; 
\ding{174} $y=x(2-x)(x\geq 0)$; 
\ding{175} $y=\begin{cases}
2x, & x\geq 2\\
4, & x=1
\end{cases}$. 
其中不存在反函数的序号是\blkda{\ding{174}\ding{175}}. 


\item 求下列函数的反函数: \\
\begin{enumerate*}[(1)]
\item $y=\dfrac{3x-1}{2x+1}(x\neq -\frac{1}{2})$; 

\item $y=\begin{cases}
x^2-1, & 0\leq x\leq 1\\
x^{\frac{2}{3}}, & -1\leq x<0
\end{cases}$; 

\item $y=x\abs{x}+2x$.
\end{enumerate*}
\begin{jieda}
\begin{enumerate}[(1)]
\item $y=\dfrac{1+x}{3-2x}$; 
\item $y=\begin{cases}
\sqrt{x+1}, & -1\leq x\leq 0\\
-x^{\frac{3}{2}}, & 0<x\leq 1
\end{cases}$; 
\item $y=\begin{cases}
1-\sqrt{1-x}, & x<0\\
\sqrt{1+x}-1, & x\geq 0
\end{cases}$. 
\end{enumerate}
\end{jieda}
\begin{vspace_}
    \vsp[3]
\end{vspace_}


\item 已知 $f(x)=x^2-2 a x+2, x \in[-1,1]$.
\begin{enumerate}
\item 当 $a=1$ 时, 函数 $y=f(x)$ 的反函数为 $y=f^{-1}(x)$, 求 $f^{-1}$ (1);
\item 若函数 $y=f(x)$ 存在反函数，求对应的$a$的取值范围以及对应的反函数 $y=f^{-1}(x)$.
\end{enumerate}
\begin{jieda}
    \begin{enumerate}
        \item 当$a = 1$时，$f(x) = (x-1)^2+1 (x \in [-1, 1])$，故其反函数为$f^{-1}(x) = -\sqrt{x-1}+1 (x \in [1, 5])$
        \item 当$a \geq 1$时，$f^{-1}(x) = -\sqrt{x+a^2-2}+a (x \in [3-2a, 3+2a])$ \\
        当$a \leq -1$时，$f^{-1}(x) = \sqrt{x+a^2-2}+a (x \in [3+2a, 3-2a])$
    \end{enumerate}
\end{jieda}
\begin{vspace_}
    \vsp[3]
\end{vspace_}


\item 设函数$f(x)=9^x-2\times3^x(x\zr)$, 该函数是否有反函数? 如果有, 求出反函数; 如果没有, 请改变其定义域, 使其有反函数. 
\begin{jieda}
$f(\log_3\frac{1}{2})=f(\log_3\frac{3}{2})=-\frac{3}{4}$, 故不存在反函数. \\
$f(x)=(3^x-1)^2-1$. \\
$(-\infty,0]$或$[0,+\infty)$的子集. 
\end{jieda}

\item 已知 $f(x)=x^2, g(x)=-\frac{1}{2} x+5, y=g^{-1}(x)$ 表示 $y=g(x)$ 的反函数, 设 $F(x)=f\left[g^{-1}(x)\right]-g^{-1}[f(x)]$, 求 $F(x)$ 的最小值.
\begin{jieda}
    $g^{-1}(x) = -2(x-5)$，$F(x) = \left[-2(x-5)\right]^2 + 2(x^2-5) = 6x^2 - 40x + 90$ \\
    故当$x = \dfrac{10}{3}$时，$F(x)$最小值为$\dfrac{70}{3}$
\end{jieda}
\begin{vspace_}
    \vsp[3]
\end{vspace_}

\item 已知 $f(x)=2^x+k$ ( $k$ 为常数), 函数 $y=f(x)$ 的反函数为 $y=f^{-1}(x)$, 且 $A(-k, 2)$ 是函数 $y=f^{-1}(x)$ 图像上的点.
\begin{enumerate}
\item 求实数 $k$ 的值及函数 $y=f^{-1}(x)$ 的解析式;
\item 将 $y=f^{-1}(x)$ 向右平移 2 个单位, 得到函数 $y=g(x)$ 的图像, 若不等式 $f^{-1}(x)-g(\sqrt{x}) \leq m$ 有解, 试求实数 $m$ 的取值范围.
\end{enumerate}
\begin{jieda}
\begin{enumerate}
    \item $k = -2$，$f^{-1}(x) = log_2(x+2)(x>-2)$ 
    \item $m \geq \dfrac{3}{2}$
\end{enumerate}
\end{jieda}
\begin{vspace_}
    \vsp[3]
\end{vspace_}

\end{enumerate}

C组：
\begin{enumerate}
    \item 全集$U$有2024个元素，若$|A| = 2000$，$|B| = 1978$，$|C| = 1958$，求$|A\cap B \cap C|_{min}$．
    \begin{jieda}
    $|A \cap B| = |A| + |B| - |A \cup B| \geq 2000 + 1978 - 2024 = 1954$ \\
    $|A \cap C| = |A| + |C| - |A \cup C| \geq 2000 + 1958 - 2024 = 1934$ \\
    $|B \cap C| = |B| + |C| - |B \cup C| \geq 1978 + 1958 - 2024 = 1912$ \\
    $|A\cap B \cap C| = |U| - |A| - |B| - |C| + |A \cap B| + |A \cap C| + |B \cap C| \geq 1888$ \\
    当$A \cup B = A \cup C = B \cup C = U$时取等号.
    \end{jieda}
    \begin{vspace_}
        \vsp[3]
    \end{vspace_}
    \item 求函数$f(x) = x^2 + \sqrt{x^4 - 3x^2+2x+5}$的值域.
    \begin{jieda}
    $f(x) = x^2 + \sqrt{(x^2-2)^2+(x+1)^2}$，其几何意义为抛物线$y = x^2$上点$B$到$A(-1, 2)$的距离与$B$到$x$轴距离之和. 由图像可以判断出$f(x) \in [2, +\infty)$ 
    
    \end{jieda}
    \item 已知$x^2+y^2+z^2 \leq 1$，求$x^2+2y-2z+3$的取值范围.
    \begin{jieda}
        $[3-2\sqrt{2}, 3+2\sqrt{2}]$
    \end{jieda}
\end{enumerate}

\section*{反函数（2）}

\zsd 反函数的图像

A组：

\begin{enumerate}
\item 若函数 $y=f(x)$ 存在反函数 $y=f^{-1}(x)$, 则函数 $y=f(x)$ 和 $y=f^{-1}(x)$\dotfill(\kh{B})
\xx{不能关于原点对称}
{单调性不可能相反}
{不可能同时是奇函数}
{如果图像存在交点, 则交点一定在 $y=x$ 直线}


\item $y=f(x)$ 有反函数 $y=f^{-1}(x)$, 且 $(a, f(a))$ 在 $y=f(x)$ 图像上, 则必在 $y=f^{-1}(x)$ 的图象上的点是\blkda{$(f(a), a)$}.

\item 已知 $f(x)=\sqrt{2 x-1}$, 函数 $y=f(x)$ 的反函数为 $y=f^{-1}(x)$, 则 $f^{-1}(3)=$ \blkda{$5$}.

\item 已知 $f(x)=1+\frac{1}{x}$, 若函数 $y=f(x)$ 的反函数为 $y=f^{-1}(x)$, 则方程 $f^{-1}(x)=2$ 的解 $x=$ \blkda{$\frac{3}{2}$}.

\item 设函数$f(x)$的图像关于点$(1,2)$对称, 且存在反函数$f^{-1}(x)$, $f(4)=0$, 则$f^{-1}(4)=$\blkda[0.7]{$-2$}. 
\begin{jieda}
$(4,0)\xleftrightarrow{\text{点}(1,2)}(-2,4)$, 故$f(-2)=4$, 于是$f^{-1}(4)=-2$. 
\end{jieda}

\end{enumerate}

B组：

\begin{enumerate}
\item 已知点$(2,\frac{1}{4})$既在函数$y=2^{ax+b}$的图像上, 又在它的反函数的图像上, 则实数$a=$\blkda[0.7]{$-\frac{12}{7}$}, $b=$\blkda[0.7]{$\frac{10}{7}$}. 

\item 已知函数$y=f(x)$存在反函数$y=f^{-1}(x)$, 若函数$y=f(1+x)$的图像经过点$(3,1)$, 则函数$y=f^{-1}(x)$的图像必经过点\blkda{$(1,4)$}. 

\item 设函数$y=f(x)$存在反函数$y=f^{-1}(x)$, 且函数$y=x-f(x)$的图像过点$(1,2)$, 则函数$y=f^{-1}(x)-x$的图像一定过点\blkda{$(-1,2)$}. 
\begin{jieda}
$2=1-f(1)\iff f(1)=-1$, 于是$f^{-1}(-1)=1$, 故$f^{-1}(-1)-(-1)=2$. 
\end{jieda}

\item 函数$f(x)=a^x-\frac{5}{2}a+6$的反函数$f^{-1}(x)$的图像经过点$(5,2)$, 且在区间$(\frac{23}{4},+\infty)$上恒有$f^{-1}(x)<0$, 则实数$a=$\blkda{$\frac{1}{2}$}. 
\begin{jieda}
$f(2)=5\iff a^2-\frac{5}{2}a+6=5\iff a=\frac{1}{2},2$. \\
当$a=\frac{1}{2}$时, $f(x)=(\frac{1}{2})^x+\frac{19}{4}$, $f^{-1}(x)=\log_{\frac{1}{2}}(x-\frac{19}{4})$, 满足题意; \\
当$a=2$时, $f(x)=2^x+1$, $f^{-1}(x)=\log_2(x-1)$, $f^{-1}(9)=3>0$, 不满足题意, 故舍去. \\
综上, $a=\frac{1}{2}$. 
\end{jieda}


\item 函数$f(x)=\sqrt{x-a}(x\geq a)$的图像与其反函数的图像有公共点, 则实数$a$的范围是\blkda{$a\leq\frac{1}{4}$}. 
\begin{jieda}
因为单调递增, 故原与反的交点在$y=x$上, 于是只要方程$\sqrt{x-a}=x(x\geq a)$有解即可. \\
$\iff t=t^2+a(t=\sqrt{x-a}\geq 0)\iff a=t-t^2(t\geq 0)$, 于是$a\leq\frac{1}{4}$. 
\end{jieda}

\item 函数$y=\dfrac{\me^x-\me^{-x}}{2}$的反函数\dotfill(\kh{C}) 
\xx{是奇函数, 在$(0,+\infty)$上是减函数}
{是偶函数, 在$(0,+\infty)$上是减函数}
{是奇函数, 在$(0,+\infty)$上是增函数}
{是偶函数, 在$(0,+\infty)$上是增函数}


\item 奇函数$y=f(x)$有反函数$y=f^{-1}(x)$, 则必在$y=f^{-1}(x)$的图像上的点是\dotfill(\kh{B})
\xx{$(-f(a),a)$}
{$(-f(a),-a)$}
{$(-a,f^{-1}(a))$}
{$(a,f^{-1}(-a))$}


\item 已知函数$f(x)=\dfrac{1-2x}{1+x}$, 函数$y=g(x)$的图像与$y=f^{-1}(x-1)$的图像关于直线$y=x$对称, 则$g(x)$等于\dotfill(\kh{B})
\xx{$\dfrac{3-2x}{x}$}
{$\dfrac{2-x}{x+1}$}
{$\dfrac{1-x}{x+2}$}
{$\dfrac{3}{x+2}$}
\begin{jieda}
$y=f^{-1}(x-1)\iff x-1=f(y)\iff x=f(y)+1$, 故$g(x)=f(x)+1$. 
\end{jieda}


\item 下面四个命题: 
\begin{dingautolist}{172}
\item 关于直线$y=x$成轴对称的两个图形一定是互为反函数的一对函数的图像;
\item 函数$y=f(x)(x\in D)$与函数$x=f^{-1}(y)(y\in f(D))$是互为反函数; 
\item 函数$y=f(x)$与其反函数$y=f^{-1}(x)$的图像若有交点, 则交点一定在直线$y=x$上;
\item 函数$y=f(x+a),x\zr$与函数$y=f^{-1}(x+a),x\zr$的图像关于直线$y=x$对称. 
\end{dingautolist}
其中错误的命题有\blkda{\ding{172}\ding{173}\ding{174}\ding{175}}.

\item 已知直线 $f(x)=k x+b$ 与曲线 $g(x)=\frac{a}{x}(a>0)$ 交于 $A\left(x_1,-3\right), B\left(x_2, 4\right)$ 两点,则不等式 $f^{-1}(x) \leq g^{-1}(x)$ 的解集为\blkda[3]{$(-\infty,-3] \cup(0,4]$}.

\item 设函数 $y=f(x)$ 存在反函数 $y=f^{-1}(x)$, 且函数 $y=x-f(x)$ 的图像过点 $(1,3)$,则函数 $y=f^{-1}(x)+3$ 的图像一定经过定点\blkda{$(-2,4)$}.


\item 给定实数$a\neq 0,a\neq 1$, 设函数$f(x)=\dfrac{x-1}{ax-1}(x\zr,x\neq\frac{1}{a})$. 求证: 函数$f(x)$的图像关于$y=x$对称. 
\begin{jieda}
$f^{-1}(x)=f(x)$. 
\end{jieda}

\item 函数$f(x)$在$\bR$上是减函数, 且它的反函数为$f^{-1}(x)$, 如果$A(-2,1)$, $B(2,-3)$是$y=f(x)$图像上的两点, 求不等式$\abs{f^{-1}(\frac{x-1}{x})}<2$的解集. 
\begin{jieda}
$\iff -2<f^{-1}(\frac{x-1}{x})<2\iff f^{-1}(1)<f^{-1}(\frac{x-1}{x})<f^{-1}(-3)
\xLongleftrightarrow{\text{单调递减}} 1>\frac{x-1}{x}>-3$\\
解得$x\in(\frac{1}{4},+\infty)$. 
\end{jieda}


\item 已知函数$f(x)=\left(\dfrac{x-1}{x+1}\right)^2(x>1)$. 
\begin{enumerate}[(1)]
\item 求函数$f(x)$的反函数$f^{-1}(x)$; 
\item 如果不等式$(1-\sqrt{x})\cdot f^{-1}(x)>m(m-\sqrt{x})$对$\left[\dfrac{1}{4},\dfrac{1}{2}\right]$上的每一个$x$都成立, 求$m$的范围;
\item 设函数$g(x)=\frac{1}{f^{-1}(x)}+\sqrt{x}+2$, 求函数$g(x)$的最小值及相应的$x$的值.
\end{enumerate}
\begin{jieda}
\begin{enumerate}[(1)]
\item $f^{-1}(x)=\dfrac{1+\sqrt{x}}{1-\sqrt{x}}(0<x<1)$. 

\item $\iff 1+\sqrt{x}>m(m-\sqrt{x})\iff(m+1)\sqrt{x}>m^2-1$. \\
于是只要$(m+1)\sqrt{\frac{1}{4}}>m^2-1\land(m+1)\sqrt{\frac{1}{2}}>m^2-1$即可, 解得$m\in(-1,\frac{3}{2})$. 

\item $g(x)=\sqrt{x}+1+\dfrac{2}{\sqrt{x}+1}$, 不难求得$g(x)_{\min}=g(3-2\sqrt{2})=2\sqrt{2}$. 
\end{enumerate}
\end{jieda}


\item 已知函数 $y=f(x), x \in D$ 是严格递增函数, 其反函数是 $y=f^{-1}(x)$.
\begin{enumerate}
\item 若 $y=x^2-1\left(x>\frac{1}{2}\right)$, 求 $y=f^{-1}(x)$, 并写出定义域 $M$;
\item 对于 (1) 的 $y=f^{-1}(x)$ 和 $M$, 设任意 $x_1 \in M, x_2 \in M, x_1 \neq x_2$, 求证:
$$
\left|f^{-1}\left(x_1\right)-f^{-1}\left(x_2\right)\right|<\left|x_1-x_2\right| .
$$
\end{enumerate}
\begin{jieda}
\begin{enumerate}
\item $f^{-1}(x)=\sqrt{x+1}, \quad M=\left(-\frac{3}{4},+\infty\right)$.
\item 证明：任意取$x_1 \in M, x_2 \in M, x_1 \neq x_2$，则
$$
\begin{aligned}
& \left|f^{-1}\left(x_1\right)-f^{-1}\left(x_2\right)\right|=\left|\sqrt{x_1+1}-\sqrt{x_2+1}\right|=\frac{\left|x_1-x_2\right|}{\sqrt{x_1+1}+\sqrt{x_2+1}} \\
\because & x_1>-\frac{3}{4}, \quad \therefore \sqrt{x_1+1}>\frac{1}{2}, \quad x_2>-\frac{3}{4}, \quad \therefore \sqrt{x_2+1}>\frac{1}{2}, \quad \therefore \sqrt{x_1+1}+\sqrt{x_2+1}>1, \quad \therefore \\
& 0<\frac{1}{\sqrt{x_1+1}+\sqrt{x_2+1}}<1 \\
\therefore & \frac{\left|x_1-x_2\right|}{\sqrt{x_1+1}+\sqrt{x_2+1}}<\left|x_1-x_2\right|, \quad \therefore\left|f^{-1}\left(x_1\right)-f^{-1}\left(x_2\right)\right|<\left|x_1-x_2\right|
\end{aligned}
$$
\end{enumerate}

\end{jieda}


\item 已知函数$f(x)=\sqrt{ax+2}(a<0)$, 其反函数为$f^{-1}(x)$. 
\begin{enumerate}[(1)]
\item 若点$P(\sqrt{3},-1)$在反函数$f^{-1}(x)$的图像上, 求$a$的值; 
\item 求证: 函数$f(x)$的图像与$y=x$的图像有且仅有一个公共点; 
\item 如果点$(m,n)(m\neq n)$是函数$f(x)=\sqrt{ax+2}(a<0)$与其反函数$f^{-1}(x)$图像上的公共点, 求$a$的取值范围.
\end{enumerate}
\begin{jieda}
\begin{enumerate}[(1)]
\item $f(-1)=\sqrt{3}$, 故$a=-1$. 

\item $y=f(x)-x=\sqrt{ax+2}-x$严格单调递减. \\
$x=0$时, $y=\sqrt{2}>0$; $x=-\frac{2}{a}$时, $y=\frac{2}{a}<0$. \\
故有根且唯一. 

\item 
    $f^{-1}(x) = \dfrac{x^2-2}{a} (x \geq 0)$，由点$(m, n) (m \neq n)$是公共点可知$f(m) = n, f^{-1}(m) = n$，根据反函数的性质可知$f(n) = m, f^{-1}(n) = m$ \\
        选取$f^{-1}(m) = n, f^{-1}(n) = m$构建方程组（方便计算），即$\left \{ \begin{array}{ll}
             \dfrac{m^2-2}{a} = n  \\
             \dfrac{n^2-2}{a} = m
        \end{array}\right. \quad (m \geq 0, n \geq 0)$
        两式相减有$\dfrac{m^2-n^2}{a} = n-m$，即$\dfrac{(m+n)(m-n)}{a} = n-m$，因此$m = n$或$m+n = -a$，$m = n$对应着原函数与反函数在$y = x$的交点，必然存在，而$m+n = -a$对应着题目所求的不在$y = x$上的交点，不一定存在（可以通过带回$\dfrac{m^2-2}{a} = n$检验判断) \\
        因为$(m, n), (n, m)$均在$f(x)$上，故可以看做是直线$y = -(x-m)+n$与$f(x)$的交点，代入$m+n = -a$即$y = -x-a$与$f(x)$有两个交点. 从图像上可以判断出两个边界分别是$y = -x-a$与$f(x)$相切或者$y = -x-a$过$f(x)$的端点$(-\dfrac{2}{a}, 0)$. \\
        前者对应着方程$ax+2 = x^2+2ax+a^2$判别式为0，即$a^2-4(a^2-2)= 0$，解得$a = -\sqrt{\dfrac{8}{3}}$；后者对应着$\dfrac{2}{a} = a$，解得$a = -\sqrt{2}$，因此$a \in (-\dfrac{2\sqrt{6}}{3}, -\sqrt{2}]$
\end{enumerate}
\end{jieda}


\end{enumerate}



\section*{补充专题：函数的图像}

\zsd 作图、识图、数形结合

\begin{enumerate}
\item 要得到函数 $y=f(x)$ 的图像, 只须将函数 $y=f(x+1)+2$ 的图像先向 \blkda{右}平移\blkda{1}个单位, 再向\blkda{下}平移\blkda{2}个单位即可.
\item 若函数 $y=x^2+(a+2) x+3, a \leq x \leq b$ 的图像关于直线 $x=1$ 对称, 则 $b=$ \blkda{6}.
\begin{jieda}
    二次函数对称轴为$x = -\dfrac{a+2}{2}$，故$1 = -\dfrac{a+2}{2}$，即$a = -4$，故$b = 6$
\end{jieda}
\item 函数 $f(x)=x^3+1$ 的图像关于点 \blkda{$(0, 1)$}对称.
\item 若函数 $y=f(x-1)$ 是奇函数, 则函数 $y=f(x)$ 的图像关于点\blkda{$(-1, 0)$}对称.

\item 函数 $y=f(2-x)$ 的图像与函数 $y=f(x-1)$ 的图像关于\blkda[3]{直线$x = \frac{3}{2}$}对称，
\item 若方程 $|x-8|=a$ 在 $[7,9]$ 上有两个实数解, 则 $a$ 的取值范围是 \blkda{$(0, 1]$}.
\item 若把函数 $y=f(x)$ 的图像作平移, 可以使图像上的点 $P(1,0)$ 变换成点 $Q(2,2)$,则函数 $y=f(x)$ 的图像经此变换后所得图像对应的函数为\dotfill(\kh{A})
\xx{$y=f(x-1)+2$}
{$y=f(x-1)-2$}
{$y=f(x+1)+2$}
{$y=f(x+1)-2$}

\item 函数$y=\dfrac{x-2}{x-1}$的图像是\dotfill(\kh{B})
\fourtu
{\includegraphics[totalheight=4.5\baselineskip]{./bodyxb/src/bx01ch03hanshudetuxiangt8a.png}}
{\includegraphics[totalheight=4.5\baselineskip]{./bodyxb/src/bx01ch03hanshudetuxiangt8b.png}}
{\includegraphics[totalheight=4.5\baselineskip]{./bodyxb/src/bx01ch03hanshudetuxiangt8c.png}}
{\includegraphics[totalheight=4.5\baselineskip]{./bodyxb/src/bx01ch03hanshudetuxiangt8d.png}}
\begin{jieda}
    可以先确定对称中心是$(1, 1)$，然后再代入$x = 0$判断.
\end{jieda}

\item 已知函数$y=f(\abs{x})$的图像如左下图所示, 则函数$y=f(x)$的图像不可能是\dotfill(\kh{B})
\begin{minipage}{\linewidth}
\centering
\includegraphics[height=5.0\baselineskip]{./bodyxb/src/bx01ch03hanshudetuxiangt9start.png}\qquad
\includegraphics[height=5.0\baselineskip]{./bodyxb/src/bx01ch03hanshudetuxiangt9a.png}\qquad
\includegraphics[height=5.0\baselineskip]{./bodyxb/src/bx01ch03hanshudetuxiangt9b.png}\qquad
\includegraphics[height=5.0\baselineskip]{./bodyxb/src/bx01ch03hanshudetuxiangt9c.png}\qquad
\includegraphics[height=5.0\baselineskip]{./bodyxb/src/bx01ch03hanshudetuxiangt9d.png}
\end{minipage}
\begin{jieda}
    $f(|x|)$是将$f(x)$在$y$轴右侧图像翻到左侧
\end{jieda}

\item 已知函数$f(x)=ax^3+bx^2+cx+d$的图像如下图, 则\dotfill(\kh{A})
\xx{$b\in(-\infty,0)$}{$b\in(0,1)$}{$b\in(1,2)$}{$b\in(2,+\infty)$}
\begin{minipage}{\linewidth}
\centering
\includegraphics[height=5\baselineskip]{./bodyxb/src/bx01ch03hanshudetuxiangt11.png}
\end{minipage}
\begin{jieda}
    首先可以判断出$a > 0$，然后根据三个零点列方程可以得到$b = -3a$或者将$f(x)$写成因式相乘的形式：$f(x) = ax(x-1)(x-2)$，也可以得到$b = -3a$
\end{jieda}

\item 任取$x_1,x_2\in[a,b], x_1\neq x_2$, 若$f\left(\dfrac{x_1+x_2}{2}\right)>\dfrac{f(x_1)+f(x_2)}{2}$, 则称$f(x)$是$[a,b]$上的凸函数. 
下列图像中, 可以是凸函数图像的是\dotfill(\kh{C})
\fourtu
{\includegraphics[totalheight=5.5\baselineskip]{./bodyxb/src/bx01ch03hanshudetuxiangt10a.png}}
{\includegraphics[totalheight=5.5\baselineskip]{./bodyxb/src/bx01ch03hanshudetuxiangt10b.png}}
{\includegraphics[totalheight=5.5\baselineskip]{./bodyxb/src/bx01ch03hanshudetuxiangt10c.png}}
{\includegraphics[totalheight=5.5\baselineskip]{./bodyxb/src/bx01ch03hanshudetuxiangt10d.png}}

\item 函数 $f(x)=2^{|x-1|}$ 的单调递减区间是\blkda{$(-\infty,1)$}.
\begin{jieda}
    先判断对称性，再分析单调性
\end{jieda}

\item 将函数 $g(x)=3^{-2 x}$ 的图象经过下列哪一种变换可以得到函数 $f(x)=3^{2-2 x}$的图象\dotfill(\kh{B})
\xx{向左平移 $1$ 个单位长度}
{向右平移 $1$ 个单位长度}
{向左平移 $2$ 个单位长度}
{向右平移 $2$ 个单位长度}

\item 已知函数 $f(x)=4^x-a \cdot 2^x+a$, 在 $x \in(0,+\infty)$ 的图像恒在 $x$ 轴上方, 则实数 $a$的取值范围是\blkda{$(-\infty,4)$}.

\item 已知函数 $f(x)=1-\frac{2}{2^x+1}$, 则\dotfill(\kh{C})
\xx{$f(x)$ 是偶函数且是增函数}
{$f(x)$ 是偶函数且是减函数}
{$f(x)$ 是奇函数且是增函数}
{$f(x)$ 是奇函数且是减函数}

\item 设函数 $f(x)=3^{-x^2+a x}$ 在区间 $(1,2)$ 上单调递增, 则 $a$ 的取值范围是\blkda{$[4,+\infty)$}.
\begin{jieda}
函数 $y=3^x$ 在 R 上单调递增, 而函数 $f(x)=3^{-x^2+a x}$ 在区间 $(1,2)$ 上单调递增,因此 $\frac{a}{2} \geq 2$, 解得 $a \geq 4$, 所以 $a$ 的取值范围是 $[4,+\infty)$. 
\end{jieda}

\item 已知 $f(x)=\left\{\begin{array}{c}4^x-2^{x+2}+m, x \leq 0 \\ x+\frac{1}{x}, x>0\end{array}\right.$ 的最小值为 2 , 则 $m$ 的取值范围为 \blkda{$[5,+\infty)$}.

\item 作出下列函数的图像:
\begin{enumerate}
\item $y=\frac{3 x+1}{x-2}$;

\begin{minipage}{\linewidth}
\centering
\includegraphics[totalheight=3cm]{./bodyxb/src/hstx18.jpg}
\end{minipage}

\item $y=|x-2|+|2 x+1|$;

\begin{minipage}{\linewidth}
\centering
\includegraphics[totalheight=3cm]{./bodyxb/src/hstx18.jpg}
\end{minipage}
\item $y=|x-2| \cdot(x+1)$

\begin{minipage}{\linewidth}
\centering
\includegraphics[totalheight=3cm]{./bodyxb/src/hstx18.jpg}
\end{minipage}

\item $y=\frac{1-|x|}{|1-x|}$

\begin{minipage}{\linewidth}
\centering
\includegraphics[totalheight=3cm]{./bodyxb/src/hstx18.jpg}
\end{minipage}

\item $y=x^2-2|x|-3$

\begin{minipage}{\linewidth}
\centering
\includegraphics[totalheight=3cm]{./bodyxb/src/hstx18.jpg}
\end{minipage}


\item $y=\left|x^2-2 x-3\right|$

\begin{minipage}{\linewidth}
\centering
\includegraphics[totalheight=3cm]{./bodyxb/src/hstx18.jpg}
\end{minipage}


\item $y=\left|x^2-2\right| x|-3|$

\begin{minipage}{\linewidth}
\centering
\includegraphics[totalheight=3cm]{./bodyxb/src/hstx18.jpg}
\end{minipage}


\end{enumerate}

\item 
\begin{enumerate}
\item 设 $[x]$ 表示不超过 $x$ 的最大整数，试作出函数 $f(x)=[x], g(x)=x-[x]$ 的图象；

\begin{minipage}{\linewidth}
\centering
\includegraphics[totalheight=3cm]{./bodyxb/src/hstx19.jpg}
\end{minipage}

\item 设 $\{x\}$ 表示不小于 $x$ 的最小整数，试作出函数 $f(x)=\{x\} , g(x)=x+\{x\}$的图象.

\begin{minipage}{\linewidth}
\centering
\includegraphics[totalheight=3cm]{./bodyxb/src/hstx19.jpg}
\end{minipage}

\end{enumerate}

\begin{jieda}
\begin{enumerate}[(1)]
\item $f(x+k)=k+f(x)(k\zz)$, $g(x+k)=x+k-f(x+k)=x-f(x)=g(x)(k\zz)$. \\
\begin{minipage}{\textwidth}
\includegraphics[height=4.5cm]{./bodyxb/src/bx01ch21jigetesuhanshu06.png}
\end{minipage}
\item $\{x+k\}=\{x\}$, 故$g(x+k)=x+k+\{x+k\}=g(x)+k$($k\zz$). $x\in[0,1)$时, $g(x)=2x$. \\
\begin{minipage}{\textwidth}
\includegraphics[height=4.7cm]{./bodyxb/src/bx01ch21jigetesuhanshu07.png}
\end{minipage}
\end{enumerate}
\end{jieda}
\begin{vspace_}
    \vsp[3]
\end{vspace_}

\item 已知函数$f(x)=x^2+2x(x\zr)$. 
\begin{enumerate}[(1)]
\item 若函数$g(x)$的图像与函数$f(x)$的图像关于原点对称, 求$g(x)$的解析式;
\item 若函数$g(x)$的图像与函数$f(x)$的图像关于$y$轴对称, 求$g(x)$的解析式;
\item 若函数$g(x)$的图像与函数$f(x)$的图像关于直线$x=1$对称, 求$g(x)$的解析式.
\end{enumerate}
\begin{jieda}
\begin{enumerate}[(1)]
\item $g(x)=-f(-x)=-x^2+2x$; 
\item $g(x)=f(-x)=x^2-2x$; 
\item $g(x)=f(2-x)=(2-x)^2+2(2-x)=x^2-6x+8$. 
\end{enumerate}
\end{jieda}
\begin{vspace_}
    \vsp[3]
\end{vspace_}

\item 已知函数 $f(x)=x^2+a x+3-a, a \in R$.
\begin{enumerate}
    \item 当 $0 \leq x \leq 2$ 时, 函数 $y=f(x)$ 的最小值是关于 $a$ 的函数 $m(a)$. 求 $m(a)$ 的最大值及其相应的 $a$ 值;
    \item 对于 $a \in R$, 研究函数 $y=f(x)$ 的图像与函数 $y=\left|x^2-2 x-3\right|$ 的图像公共点的个数; (写出结论即可).
\end{enumerate}
\begin{jieda}
    \begin{enumerate}
        \item 二次函数的对称轴为$x = -\dfrac{a}{2}$，故$m(a) = \left \{ \begin{array}{ll}
            f(0), & -\dfrac{a}{2} < 0 \\
            f( -\dfrac{a}{2}), & 0 \leq -\dfrac{a}{2} \leq 2 \\
            f(2), & -\dfrac{a}{2} > 2
        \end{array}\right.$，\\
        即$m(a) = \left \{ \begin{array}{ll}
            3-a, & a > 0 \\
            -\dfrac{a^2}{4}-a+3, & -4 \leq a \leq 0 \\
            7+a, & a < -4
        \end{array}\right.$，故$m(a)$最大值为$4$，在$a = -2$处取到.
        \item 当$a = -2$时，有$1$个交点，当$a \in (-6, -2)\cup(-2, 2)$时，有$3$个交点；当$a = 2$或$a = -6$时有$2$个交点，当$a > 2$或$a < -6$时有$1$个交点.
        % 设$g(x) = \left|x^2-2 x-3\right|$，两个函数都会过定点$(1, 4)$. \\
        % 当$a = -2$时两函数对称轴均为$x = 1$，两侧完全对称，考虑在$x = 1$右侧，$f(x)$与$g(x)$在相同位置其变化率的绝对值相同，而$g(x)$先减后增，因此不会有交点. \\
        % 当$a \in (-6, 2)$且$a \neq -2$时，$f(x)$对称轴在$(-1, 3)$之间，且顶点在$x$轴上方，考虑$a > -2$的场景，在$f(x)$对称轴右侧，$f(3) > g(3)$，根据零点存在定理可知$(-\dfrac{a}{2}, 3)$之间图像会有一个交点。当$x > 3$，在相同位置$g(x)$增长更快，故会有一个横坐标大于$3$的交点，故此时有$3$个交点. \\
        % 当$a \in \{-6, 2\}$时，$f(x)$的顶点位于$x$轴与$g(x)$的交点。考虑$a > -2$的场景，$f(3) = g(3)$，当$x > 3$，在相同位置$g(x)$增长更快，故只有两个交点。
    \end{enumerate}
\end{jieda}
\begin{vspace_}
    \vsp[3]
\end{vspace_}

\item (附加)设 $D$ 是含数 1 的有限实数集， $f(x)$ 是定义在 $D$ 上的函数.若 $f(x)$ 的图像绕原点逆时针旋转 $\frac{\pi}{6}$ 后与原图像重合, 则在以下各项中, $f(1)$ 的可能取值是\dotfill(\kh{B})
\xx{$\sqrt{3}$}
{$\frac{\sqrt{3}}{2}$}
{$\frac{\sqrt{3}}{3}$}
{0}
\begin{jieda}
    函数图象应该是在以原点为圆心的圆上选取的12个点.\\
    设点$A(1, f(1))$，则射线$OA$与$x$轴正方向的夹角不可以等于$\dfrac{k\pi}{6}(k \in \mathbf{Z})$，因此选B
\end{jieda}


\item (附加) 设函数 $f_0(x)=\left(\frac{1}{2}\right)^{|x|}, f_1(x)=\left|f_0(x)-\frac{1}{2}\right|, f_n(x)=\left|f_{n-1}(x)-\left(\frac{1}{2}\right)^n\right|, n \geq 1, n \in N$, 则方程 $f_n(x)=\left(\frac{1}{n+2}\right)^n$ 有\blkda{$2^{n+1}$}个实数根。


\end{enumerate}

\section*{对数定义与运算}

\zsd 对数、对数的运算(换底公式）

A组：

\begin{enumerate}
\item 若$2^x=5$, 则$x$=\blkda{$\log_25$}; 若$\log_2x=4$, 则$x$=\blkda{16}. 

\item 若$\log_x3=\dfrac{1}{4}$, 则$x$=\blkda{81}; 若$\log_2x=-4$, 则$x$=\blkda{$\dfrac{1}{16}$}. 

\item 若$\log_7x=0$, 则$x$=\blkda{1}; 若$\log_x100=1$, 则$x$=\blkda{100}. 

\item 计算: $\log_2\dfrac{1}{64}$=\blkda{$-6$}; $\log_{\sqrt{3}}\frac{\sqrt{3}}{3}$=\blkda{$-1$}. 

\item 计算: $3^{\log_317}$=\blkda{17}. 


\end{enumerate}


B组：

\begin{enumerate}
\item 若$\log_ac\cdot\log_bc=0$, 则$ab+c-abc$=\blkda{1}. 
\begin{jieda}
$\log_ac=0$或$\log_bc=0$, 故$c=1$. 
\end{jieda}


\item 若$\lg{a},\lg{b}$是方程$2x^2-4x+1=0$的两个实根, 则$\lg(ab)\cdot\lg^2\frac{a}{b}$的值为\blkda{4}. 
\begin{jieda}
$\begin{cases}
\lg{a}+\lg{b}=2\\
\lg{a}\lg{b}=\frac{1}{2}
\end{cases}$, 
故$\lg(ab)\cdot\lg^2\frac{a}{b}=(\lg{a}+\lg{b})(\lg{a}-\lg{b})^2=2\cdot(2^2-4\cdot\frac{1}{2})=4$
\end{jieda}

\item $3^{\log_{\sqrt{3}}4}$的值是\blkda{$16$}.

\item $\sqrt{5}^{\log_5(-a)^2}(a\neq 0)$化简的结果是\blkda{$\abs{a}(a \neq 0)$}.


\item 若$\log_4(\log_3(\log_2x))=0$, 则$x^{-\frac{1}{2}}$等于\blkda{$\dfrac{1}{4}\sqrt{2}$}.
\begin{jieda}
$x=2^{3^{4^0}}=8$. 
\end{jieda}

\item $\log_{\sqrt{n+1}+\sqrt{n}}\left(\sqrt{n+1}-\sqrt{n}\right)$等于\blkda{$-1$}.
\begin{jieda}
设所求为$t$, 则$(\sqrt{n+1}+\sqrt{n})^t=\sqrt{n+1}-\sqrt{n}=(\sqrt{n+1}-\sqrt{n})^{-1}$, 故$t=-1$. 
\end{jieda}

\item 若$2\ln\dfrac{a-b}{2}=\ln{a}+\ln{b}$, 则$\dfrac{a}{b}$=\blkda{$3+2\sqrt{2}$}. 


\item 若$\log_427=m$, $\log_325=n$, 则$\lg2$可用$m,n$表示为\blkda{$\dfrac{3}{3+mn}$}. 
\begin{jieda}
乱凑一通: $mn=\log_427\cdot\log_325=\dfrac{\lg27}{\lg4}\cdot\dfrac{\lg25}{\lg3}=\dfrac{3\lg3\cdot2\lg5}{2\lg2\cdot\lg3}$, \\
故$\dfrac{mn}{3}=\dfrac{\lg5}{\lg2}$, 于是$1+\dfrac{mn}{3}=\dfrac{\lg5}{\lg2}+1=\dfrac{\lg5+\lg2}{\lg2}=\dfrac{1}{\lg2}$, \\
则$\lg2=\dfrac{3}{3+mn}$. 
\end{jieda}


\item 若$3^a=4^b=36$, 则$\dfrac{2}{a}+\dfrac{1}{b}$的值为\blkda{$1$}. 
\begin{jieda}
$a=\log_336$, $b=\log_436$, 则$\dfrac{2}{a}+\dfrac{1}{b}=\dfrac{2}{\log_336}+\dfrac{1}{\log_436}=2\log_{36}3+\log_{36}4=\log_{36}36=1$. 
\end{jieda}


\item 适合$\log_5x\cdot\log_x7=\log_57$的$x$的集合是\dotfill(\kh{C})
\xx{$\{5,7\}$}{$\{0,1\text{以外的实数}\}$}{$\{\text{不为}1\text{的正数}\}$}{$\bR$}


\item 设$m=\log_ax$, $n=\log_ay$, $p=\log_az$, 用$m,n,p$表示下列各式: \\
\begin{enumerate*}[(1)]
\item $\log_a\dfrac{x^2\cdot\sqrt{y}}{\sqrt[3]{z}}$; 
\item $\log_a\dfrac{az}{x^{\frac{3}{2}}y}$. 
\end{enumerate*}
\begin{jieda}
$x=a^m,y=a^n,z=a^p$. 
\begin{enumerate}[(1)]
\item 原式$=\log_a\dfrac{a^{2m}\cdot a^{\frac{n}{2}}}{a^{\frac{p}{3}}}=\log_aa^{2m+\frac{n}{2}-\frac{p}{3}}=2m+\frac{n}{2}-\frac{p}{3}$. 
\item 同理可得：原式$=1+p-\frac{3m}{2}-n$. 
\end{enumerate}
\end{jieda}

\begin{vspace_}
    \vsp[3]
\end{vspace_}


\item 已知$\log_a2=m$, $\log_a3=n$, 求$a^{2m-3n}$的值. 
\begin{jieda}
由条件可得 $a^m=2, a^n=3$, 其中$a>0\land a\neq 1$. \\
于是$a^{2m-3n}=\dfrac{2^2}{3^3}=\dfrac{4}{27}$. 
\end{jieda}
\begin{vspace_}
    \vsp[3]
\end{vspace_}

\item 已知$3^{2x}=4^{3y}=12^6$, 求$\dfrac{3}{x}+\dfrac{2}{y}$的值. 
\begin{jieda}
直接弄倒数, 少写几步过程. \\
原式$=\dfrac{6}{2x}+\dfrac{6}{3y}=6\cdot\log_{12^6}3+6\cdot\log_{12^6}4=\log_{12}3+\log_{12}4=1$. 
\end{jieda}

\begin{vspace_}
    \vsp[3]
\end{vspace_}


\item 
\begin{enumerate}
\item 若$\log_627=a$, 用$a$表示$\log_{18}16$. 
\begin{jieda}
$\dfrac{1}{a}=\log_{27}6$, 于是$\dfrac{3}{a}=1+\log_32$, 故$\log_23=\dfrac{a}{3-a}$. \\
于是$\log_{18}16=4\log_{18}2=\dfrac{4}{\log_218}=\dfrac{4}{1+2\log_23}=\dfrac{4}{1+2\cdot\frac{a}{3-a}}=\dfrac{12-4a}{3+a}$. 
\end{jieda}

\begin{vspace_}
    \vsp[3]
\end{vspace_}


\item 若$\log_{15}2=a$, $3^b=5$, 试用$a,b$表示$\log_{125}18$. 
\begin{jieda}
$a=\dfrac{\log_52}{\log_515}=\dfrac{\log_52}{\log_53+1}$, $\dfrac{1}{b}=\log_53$. \\
于是$\log_{125}18=\dfrac{1}{3}(\log_52+2\log_53)=\dfrac{1}{3}\left(a(1+\frac{1}{b})+2\cdot\frac{1}{b}\right)=\dfrac{ab+a+2}{3b}$. 
\end{jieda}
\begin{vspace_}
    \vsp[3]
\end{vspace_}
\end{enumerate}

\end{enumerate}

C组：
\begin{enumerate}
    \item 已知$log_2(x+y) = log_2x+log_2y$ \\
    \begin{enumerate*}[(1)]
        \item 求$x+2y$的最小值
        \item 求$\dfrac{4x}{x-1} + \dfrac{9y}{y-1}$的最小值
        \item 正数$z$满足$\dfrac{1}{x+y} + \dfrac{1}{z} = 1$，求$z$的取值范围
    \end{enumerate*}
    \begin{jieda}
        $log_2(x+y) = log_2(xy)$，故$x+y = xy$，即$x = \dfrac{y}{y-1}$，$y = \dfrac{x}{x-1}$，其中$x > 1, y > 1$.
        \begin{enumerate}[(1)]
            \item $x + 2y = x + \dfrac{2x}{x-1} = (x-1) + \dfrac{2}{x-1}+3 \geq 2\sqrt{2} + 3$，在$x = 1+\sqrt{2}, y = 1 + \dfrac{\sqrt{2}}{2}$时取等号.
            \item $\dfrac{4x}{x-1} + \dfrac{9y}{y-1} = \dfrac{4x}{x-1} + 9x = 9(x-1) + \dfrac{4}{x-1}+13 \geq 25$，在$x = \dfrac{5}{3}, y = \dfrac{5}{2}$时取等号.
            \item $\dfrac{1}{z} = 1 - \dfrac{1}{x+y} = 1 - \dfrac{x-1}{x^2} = 1 - \dfrac{1}{x} + \dfrac{1}{x^2}$\\
            令$t = \dfrac{1}{x}, t \in (0, 1)$，因此$\dfrac{1}{z} \in [\dfrac{3}{4}, 1)$，故$z \in (1, \dfrac{4}{3}]$
        \end{enumerate}
    \end{jieda}
    \item 已知$log_2[log_{\frac{1}{2}}(log_2x)] = log_3[log_{\frac{1}{3}}(log_3x)] = log_5[log_{\frac{1}{5}}(log_5x)] = 0$，比较实数$x, y, z$之间的大小关系.
    \begin{jieda}
        方程$log_a[log_{\frac{1}{a}}(log_ax)] = 0$的解为$a^{\frac{1}{a}}$，所以对于本题中的$x, y, z$分别为$x = \sqrt{2}, y = \sqrt[3]{3}, z = \sqrt[5]{5}$；
        故可得$y > x > z$
    \end{jieda}
    \item 解方程组$\left \{ \begin{array}{l}
         \dfrac{e}{4} = x^y \cdot lnx\\
         \dfrac{ey+eln4}{4} = x^y 
    \end{array}\right.$ \\
    \begin{jieda}
        上式除以下式可得$\dfrac{1}{y+ln4} = lnx$ \\
        为了消元，需要表示出$x^y$，故前式变形为$y+ln4 = log_xe$，进一步变形为$x^{y+ln4} = e$，因此$x^y = \dfrac{e}{x^{ln4}}$，将其代入原始方程组第一个式子有$\dfrac{e}{4} = \dfrac{e}{x^{ln4}} \cdot lnx$，即$x^{ln4} = 4lnx$，即$4^{lnx} = 4lnx$.\\
        令$t = lnx, t \in \mathbf{R}$，故$4^t = 4t$，解得$t = 1$或$2$. \\
        故$x = e, y = 1 - ln4$或$x = \sqrt{e}, y = 2 - ln4$
    \end{jieda}
    \item 已知$a_n = log_{n+1}(n+2), n \in N^*$，定义使得$a_1a_2a_3...a_k$为整数的$k$称作“希望数”,求区间$[1, 2016]$内的所有“希望数”的和.
    \begin{jieda}
        易知$a_1a_2a_3...a_k = log_2(k+2)$，因此当$k = 2^m - 2$($m \in \mathbf{N}$且$m > 1$)的时候$a_1a_2a_3...a_k$为整数.
        故$[1, 2016]$内的所有“希望数”的和为$\sum_{i = 2}^{10}2^{i} - 2 \times 9 = 2026$
    \end{jieda}
    
\end{enumerate}

\section*{对数函数}

\zsd 对数函数的定义、图像、求定义域、奇偶性、单调性

A组：
\begin{enumerate}
\item 对数函数 $f(x)=\log _3 x$ 定义域为\blkda{$(0,+\infty)$}.


\item 已知$f(x)=\begin{cases}
\log_2x, & x>0\\
3^x, & x\leq 0
\end{cases}$, 则$f(f(\frac{1}{4}))=$\blkda{$\frac{1}{9}$}. 
\begin{jieda}
$f(\frac{1}{4})=\log_2\frac{1}{4}=-2$, 则$f(f(\frac{1}{4}))=f(-2)=\frac{1}{9}$. 
\end{jieda}


\item 函数$y=\log_2(x+1)$的图像不经过第\blkda{II、IV} 象限. 
\begin{jieda}
向左平移1个单位, 图像过原点, 作图易得不过第II、IV象限. 
\end{jieda}

\item 若$f(x)=\log_3x$, 且$f^{-1}(m)=81$, 则$m=$\blkda{4}. 
\begin{jieda}
$m=f(81)=4$. 
\end{jieda}

\item 已知函数 $f(x)=\log _a x$ 在定义域内是严格减函数, 则实数 $a$ 的取值范围\blkda{$(0, 1)$}.

\item 函数 $f(x)=\log _a(2 x+1)+2(a>0$ 且 $a \neq 1)$ 必过定点\blkda{$(0, 2)$}.

\item 已知函数 $f(x)=\log _2 \frac{a-x}{1+x}$ 图像关于原点对称, 则实数 $a$ 的值为\blkda{$1$}.

\item 函数$y=\log_2x$与$y=\lg{x}$的图像的公共点是\blkda{1个}. 
\begin{jieda}
当然只有一个$(1,0)$. 
\end{jieda}


\item 若函数 $y=\log _a(x+b)(a>0, a \neq 1)$ 的图象过两点 $(-1,0)$ 和 $(0,1)$, 则\dotfill(\kh{$A$})
\xx{$a=2, b=2$}{$a=\sqrt{2}, b=2$}{$a=2, b=1$}{$a=\sqrt{2}, b=\sqrt{2}$}

\item 判断奇偶性: 函数$f(x)=\log_2\dfrac{2-x}{2+x}$是\blkda{奇} 函数.

\item 已知函数 $f(x)=\log _a x(a>0, a \neq 1)$ 的 图像过点 $(9,2)$, 则 满 足 $f(2 m-1)-f(m+3)<0$ 的实数 $m$ 的取值范围是\blkda{$(\frac{1}{2}, 4)$}.

\item 不等式 $\left|\log _2 x-a\right|<5$ 对任意 $x \in[4,16]$ 恒成立, 则实数 $a$ 的取值范围为\blkda{$(-1, 7)$}.

\end{enumerate}

B组：
\begin{enumerate}
\item 函数$f(x)=\frac{\sqrt{x-2}}{x-3}\cdot\lg\sqrt{4-x}$的定义域是\blkda[3]{$[2,3)\cup(3,4)$}. 
\begin{jieda}
$x-2\geq 0\land x-3\neq 0\land \sqrt{4-x}>0$, 故$x\in[2,3)\cup(3,4)$. 
\end{jieda}


\item 设函数$y=\lg(x^2-5x)$的定义域为$M$, 函数$y=\lg(x-5)+\lg{x}$的定义域为$N$, 则\dotfill(\kh{C})
\xx{$M\cup N=\bR$}
{$M=N$}
{$M\supseteq N$}
{$M\subseteq N$}
\begin{jieda}
$ab>0$与$\begin{cases}
a>0\\
b>0
\end{cases}$的关系. 
\end{jieda}


\item 下列各组函数中表示同一函数的是\dotfill(\kh{D})
\xx{$y=\lg10^x$, $y=10^{\lg{x}}$}
{$y=\lg2x$, $y=2\lg{x}$}
{$y=\lg{x^2}$, $y=2\lg{x}$}
{$y=\lg\sqrt{x}$, $y=\frac{1}{2}\lg{x}$}
\begin{jieda}
考虑定义域. 
\end{jieda}


\item 函数$y=\sqrt{\log_{\frac{1}{3}}\left(\log_{\frac{1}{3}}\left(\log_{\frac{1}{3}}x\right)\right)}$的定义域是\blkda{$\left(\frac{1}{3},\sqrt[3]{\frac{1}{3}}\right]$}.
\begin{jieda}
$\log_{\frac{1}{3}}\left(\log_{\frac{1}{3}}\left(\log_{\frac{1}{3}}x\right)\right)\geq 0\iff 0<\log_{\frac{1}{3}}\left(\log_{\frac{1}{3}}x\right)\leq 1
\iff 1>\log_{\frac{1}{3}}x\geq\dfrac{1}{3}\iff (\frac{1}{3})^1<x\leq(\frac{1}{3})^{\frac{1}{3}}$. 
\end{jieda}


\item 
\begin{enumerate}[(1)]
\item 函数$y=\log_2\left(x^2-2x\right)$的单调递增区间是\blkda{$(2,+\infty)$}; 
\item 函数$y=\log_{\frac{1}{2}}(x^2-5x+6)$的单调增区间为\blkda{$(-\infty,2)$}. 
\end{enumerate}


\item 已知$1<x<2$, 则$a=2^x$, $b=\sqrt[3]{x}$, $c=\log_{\frac{1}{2}}x$的大小关系是\blkda{$a>b>c$}. 


\item 函数$f(x)=\log_2(1-x)(x<-3)$的反函数是\blkda[4]{$f^{-1}(x)=1-2^x(x>2)$}.  
\begin{jieda}
$1-x>4$, 故$f(x)\in(2,+\infty)$. \\
$f^{-1}(x)=1-2^x(x>2)$. 
\end{jieda}


\item 设函数$f(x)=\log_3(x^2+1)-1(x>0)$, 求$f^{-1}(x)$=\blkda[3.5]{$\sqrt{3^{x+1}-1}(x>-1)$}.
\begin{jieda}
$f^{-1}(x)=\sqrt{3^{x+1}-1}(x>-1)$. 
\end{jieda}


\item 方程$\log_2 x=x^2-4$的解的个数是\blkda{2}个. 


\item 已知$f(x)$为奇函数, 当$-1<x<0$时, $f(x)=x\lg(x+1)$, 则当$0<x<1$时, $f(x)=$\blkda[2]{$x\lg(1-x)$}. 
\begin{jieda}
当$x\in(0,1)$时, $-x\in(-1,0)$, 于是$f(x)=-f(-x)=-(-x)\lg(-x+1)=x\lg(-x+1)$. 
\end{jieda}


\item 已知函数$f(x)=\begin{cases}
(\frac{1}{2})^x, & x\geq 4\\ 
f(x+1), & x<4\\ 
\end{cases}$, 则$f(2+\log_23)$的值为\blkda{$\frac{1}{24}$}. 
\begin{jieda}
$2+\log_23\in(3,4)$, $3+\log_23\in(4,5)$, \\
故$f(2+\log_23)=f(3+\log_23)=(\frac{1}{2})^{3+\log_23}=\frac{1}{24}$. 
\end{jieda}


\item 
\begin{enumerate}
    \item 不等式$\log_x3<\log_x2$的解集为\blkda{$(0,1)$}. 
    \item 不等式 $\lg (x-1)^2>2$ 的解集是\blkda[4]{$(-\infty, -9) \cup (1, +\infty)$}.
\end{enumerate}
\begin{jieda}
\begin{enumerate}
    \item $\iff\dfrac{\lg3}{\lg{x}}<\dfrac{\lg2}{\lg{x}}\iff\dfrac{\lg3-\lg2}{\lg{x}}<0\iff\lg{x}<0\iff 0<x<1$. 
    \item $\iff (x-1)^2 > 100 \iff x \in (-\infty, -9) \cup (1, +\infty)$. 
\end{enumerate}
\end{jieda}


\item 已知函数$g(x)=f(x)-\frac{1}{f(x)}$, 其中$\log_2f(x)=2x$, $x\zr$, 则$g(x)$\dotfill(\kh{C})
\xx{是奇函数又是减函数}{是偶函数又是增函数}{是奇函数又是增函数}{是偶函数又是减函数}
\begin{jieda}
$f(x)=4^x$, 则$g(x)=4^x-4^{-x}$. 
\end{jieda}


\item 已知函数$f(x)=\log_a(2^x+b-1)(a>0, a\neq 1)$的图像如图所示, 则$a,b$满足的关系是\dotfill(\kh{A})
\begin{minipage}{\linewidth}
\centering
\includegraphics[totalheight=2.5cm]{./bodyxb/src/main42chap21duishuhanshu01.png}
\end{minipage}
\xx{$0<a^{-1}<b<1$}
{$0<b<a^{-1}<1$}
{$0<b^{-1}<a<1$}
{$0<a^{-1}<b^{-1}<1$}
\begin{jieda}
由复合函数的单调性可知$a>1$. \\
又$-1<f(0)<0$, 故$-1<\log_a(2^0+b-1)<0$, 即$a^{-1}<b<1$. 
\end{jieda}


\item 若$\log_a2<\log_b2<0$, 则\dotfill(\kh{B})
\xx{$0<a<b<1$}
{$0<b<a<1$}
{$a>b>1$}
{$b>a>1$}
\begin{jieda}
$\iff\dfrac{1}{\log_2a}<\dfrac{1}{\log_2b}<0\iff 0>\log_2a>\log_2b\iff 1>a>b>0$. 
\end{jieda}



\item 函数$f(x)=\lg(2^x-b)$($b$为常数), 若$x\in[1,+\infty)$时, $f(x)\geq 0$恒成立, 则\dotfill(\kh{A})
\xx{$b\leq 1$}{$b<1$}{$b\geq 1$}{$b=1$}
\begin{jieda}
$f(x)$在定义域上严格单调递增, 于是只要$f(1)\geq 0$即可, 易得$b\leq 1$. 
\end{jieda}



\item 
\begin{enumerate}
\item 若函数 $f(x)=\lg \left(a x^2-4 x+a-3\right)$ 的定义域为 $R$, 则实数 $a$ 的取值范围是\blkda{$a > 4$}.
\item 若函数 $f(x)=\lg \left(a x^2-4 x+a-3\right)$ 的值域为 $R$, 则实数 $a$ 的取值范围是\blkda{$0 \leq a \leq 4$}.
\item 若函数 $y=\log _a\left(x^2-a x+1\right)$ 有最小值, 则实数 $a$ 的取值范围是\blkda{$1 < a < 2$}.
\end{enumerate}
\begin{jieda}
    \begin{enumerate}
        \item $a > 0$且$\Delta < 0$
        \item $a = 0$或$a > 0$且$\Delta \geq 0$
        \item 由复合函数单调性知$a > 1$，由最小值在内层抛物线顶点取，因此顶点应在$x$轴上方，故$\Delta < 0$
    \end{enumerate}
\end{jieda}

\item 设定义域为$\bR$的函数$f(x)=\begin{cases}
\abs{\lg\abs{x-1}}, & x\neq 1\\
0, & x=1
\end{cases}$, 则关于$x$的方程$f^2(x)+bf(x)+c=0$有7个不同实数根的充要条件是 \dotfill(\kh{C})
\xx{$b<0$且$c>0$}
{$b>0$且$c<0$}
{$b<0$且$c=0$}
{$b\geq 0$且$c=0$}
\begin{jieda}
作$y=f(x)$草图如下: \\
\begin{minipage}{\linewidth}
\centering
\includegraphics[totalheight=3cm]{./bodyxb/src/main42chap21zhiduishufangcheng01.png}
\end{minipage}
结合图像对于关于$x$的方程$t=f(x)$解的情形如下: \\
$t<0$时, 无解; $t=0$时, 3解; $t>0$时, 4解. \\
由于$3+4=7$, 于是对于$t^2+bt+c=0$而言, $t_1=0,t_2>0$. \\
故$c=0$, $-b>0$, 即$b<0,c=0$. 
\end{jieda}


\item
\begin{enumerate}[(1)]
\item 求函数$y=\dfrac{1}{\log_2(-x^2+4x-3)}$的定义域; 
\item 求函数$y=\sqrt{\log_{\frac{1}{2}}(x^2-1)}$的定义域. 
\end{enumerate}
\begin{jieda}
\begin{enumerate}[(1)]
\item $\log_2(-x^2+4x-3)\neq 0\iff x\in(1,2)\cup(2,3)$. 
\item $\log_{\frac{1}{2}}(x^2-1)\geq  0\iff 0<x^2-1\leq 1\iff x\in[-\sqrt{2},-1)\cup(1,\sqrt{2}]$. 
\end{enumerate}
\end{jieda}
\begin{vspace_}
    \vsp[3]
\end{vspace_}


\item 判断下列函数的奇偶性: 
\begin{enumerate}
\item $f(x)=\log_2(x+1)+\log_2(x-1)$; 
\item $f(x)=\lg\frac{x-1}{x+1}$;
\item 数$f(x)=\lg(\sqrt{x^2+1}-x)$.
\end{enumerate}
\begin{jieda}
\begin{enumerate}
\item 定义域不关于原点对称, 故非奇非偶. 
\item 定义域为$(-\infty,-1)\cup(1,+\infty)$, 奇函数. 
\item 定义域为$\bR$, $f(-x)=\lg(\sqrt{x^2+1}+x)=\lg\dfrac{1}{\sqrt{x^2+1}-x}=-f(x)$, 故为奇函数. 
\end{enumerate}
\end{jieda}
\begin{vspace_}
    \vsp[3]
\end{vspace_}


\item 已知$f(x)=1+\log_2x(1\leq x\leq 4)$, 求函数$g(x)=f^2(x)+f(x^2)$的最大值和最小值. 
\begin{jieda}
$x\in[1,4]\land x^2\in[1,4]$, 则$x\in[1,2]$. 令$t=\log_2x\in[0,1]$. \\
$g(x)=(1+t)^2+(1+2t)=t^2+4t+2\in[2,7]$, 故$\min=2(x=1)$, $\max=7(x=2)$. 
\end{jieda}
\begin{vspace_}
    \vsp[3]
\end{vspace_}

\item 已知$x$满足不等式$2(\log_{\frac{1}{2}}x)^2+7\log_{\frac{1}{2}}x+3\leq 0$, 
求函数$f(x)=\log_2\frac{x}{2}\log_2\frac{x}{4}$的最大值与最小值. \label{item:secduishuhanshudexingzhi13}
\begin{jieda}
令$t=\log_2x$, 则$2t^2-7t+3\leq 0$, 解得$t\in[\frac{1}{2},3]$. \\
故$y=(t-1)(t-2)=t^2-3t+2=(t-\frac{3}{2})^2-\frac{1}{4}\in[-\frac{1}{4},2]$. \\
所以$f(x)_{\min}=f(2^{\frac{3}{2}})=-\frac{1}{4}$, $f(x)_{\max}=f(8)=2$. 
\end{jieda}
\begin{vspace_}
    \vsp[3]
\end{vspace_}

\item
\begin{enumerate}[(1)]
\item 已知函数$y=\log_ax$在$x\in[2,+\infty)$时恒有$\abs{y}>1$, 求$a$的取值范围;
\item 设函数$f(x)=\lg(1+2^x+a\cdot 4^x)$, 其中$a\zr$, 如果当$x\in(-\infty,1]$时, $f(x)$恒有意义, 求$a$的取值范围.
\end{enumerate}
\begin{jieda}
\begin{enumerate}[(1)]
\item $\abs{y}=1\iff a = 2 \mbox{或} \frac{1}{2}$. \\
对数函数$y = log_ax$的图像过$(a, 1)$，在$x = 1$右侧，$a > 1$时$a$越大图像增长越慢，$1 > a > 0$时$\frac{1}{a}$越大图像下降越慢，因此由图像性质$\abs{y}>1\iff a \in (\frac{1}{2}, 1) \cup (1, 2)$

\item $1+2^x+a\cdot 4^x>0\iff a>\dfrac{-1-2^x}{4^x}=-t^2-t$, 其中$t=(\frac{1}{2})^x\in[\frac{1}{2},+\infty)$. \\
而$y=-t^2-t$在$[\frac{1}{2},+\infty)$上的最大值为$-\frac{3}{4}$(当且仅当$t=\frac{1}{2}$时取到), 故$a>-\frac{3}{4}$. 
\end{enumerate}
\end{jieda}
\begin{vspace_}
    \vsp[3]
\end{vspace_}

\item 已知$f(x)=\dfrac{1}{1+\log_2\abs{x}}+p(p\zr)$. 
\begin{enumerate}[(1)]
\item 求$f(x)$的定义域; 
\item 当$x\in(-\frac{1}{2},0)$时, 判断$f(x)$的单调性; 
\item 当$x>0$时, 若$f(x)$的反函数为$f^{-1}(x)$且$f^{-1}(0)$的值在$[2,3]$之间, 求$p$的取值范围. 
\end{enumerate}
\begin{jieda}
\begin{enumerate}[(1)]
\item $1+\log_2\abs{x}\neq 0\iff \abs{x}\neq\frac{1}{2} \mbox{且} x \neq 0 \iff x\neq\pm\frac{1}{2} \mbox{且} x \neq 0 $. 故定义域为$\{x | x \neq \pm \frac{1}{2}, x \neq 0\}$
\item 单调递增. 
\item 令$t=f^{-1}(0)$, 则$f(t)=0$, 且$t\in[2,3]$. \\
于是$\frac{1}{1+\log_2\abs{t}}\in[\log_62,\frac{1}{2}]$, 而$p=-\frac{1}{1+\log_2\abs{t}}$, 于是$p\in[-\frac{1}{2},-\log_62]$. 
\end{enumerate}
\end{jieda}
\begin{vspace_}
    \vsp[3]
\end{vspace_}

\end{enumerate}
    
\section*{指对方程}
\begin{enumerate}
    \item 方程$2^{2+x}-2^{2-x}=15$的解集是\blkda{$\{2\}$}. 
    \item 方程$9^x+6^x=2^{2x+1}$的解集是\blkda{$\{0\}$}. 
    \item 方程$2^{\left| {x+1} \right|}=3$的解集是\blkda[3]{$\left\{\log_2\frac{3}{2},\log_2\frac{1}{6}\right\}$}.
    \begin{jieda}
    $y = 2^{\left| {x+1} \right|}$ 关于$x = -1$对称，两根之和为$-2$
    \end{jieda}
    \item 方程$log_2(9-2^x) = 3-x$的解集是\blkda{$\{0, 3\}$}.
    \begin{jieda}
        转化为$9-2^x = 2^{3-x}$，即$9 = 2^x + 8 \cdot 2^{-x}$，故$2^x = 1$或$8$，$x \in \{0, 3\}$
    \end{jieda}
    \item 满足方程$3^{2x+2}-3^{x+3}-3^x+3=0$的实数解有\blkda{$2$}个.
    \item 已知$f(x)=\dfrac{1-2^x}{1+2^x}$, 则$f^{-1}(\frac{3}{5})$的值为\blkda{$-2$}.
    \begin{jieda}
    $t=f^{-1}(\frac{3}{5})\iff f(t)=\frac{3}{5}$, 解得$t=-2$. 
    \end{jieda}
    \item 求方程$9^x+4^x=\frac{5}{2}\cdot 6^x$的解集. 
    \begin{jieda}
    两边同除以$4^x$, 得$((\frac{3}{2})^x)^2-\frac{5}{2}\cdot(\frac{3}{2})^x+1=0$. \\
    于是$(\frac{3}{2})^x=2,\frac{1}{2}$, 故$x=\pm\log_{\frac{3}{2}}2$. 
    \end{jieda}
    \item 方程$4^x+4^{-x}-5(2^x+2^{-x})+6=0$的解集是\blkda[6]{$\{\log_2(2+\sqrt{3}), \log_2(2-\sqrt{3})\}$}.
    \begin{jieda}
    令$t=2^x+2^{-x}$, 则$t^2-5t+4=0$, 故$t=1$或4. \\
    而$2^x+2^{-x}=1$无解, 解$2^x+2^{-x}=4$得$x=\log_2(2\pm\sqrt{3})$. 
    \end{jieda}
    \item 关于$x$的方程$5^{-x}=\dfrac{a+3}{5-a}$的根为正数, 则实数$a$的取值范围是\blkda{$(-3,1)$}. 
    \begin{jieda}
    $x>0$时, $5^{-x}\in(0,1)$, 故$\dfrac{a+3}{5-a}\in(0,1)$, 解得$a\in(-3,1)$
    \end{jieda}
    
    \item 函数$y=3^{-\abs{x-1}}-m$的图像与$x$轴有交点时, $m$的取值范围是\blkda{$(0,1]$}. 
    \begin{jieda}
    $3^{-\abs{x-1}}\in(0,1]$, 故$m\in(0,1]$. 
    \end{jieda}
    \item 关于$x$的方程$4^x+(m+1)\cdot2^x+m=0$只有一个根, 则$m$的取值范围是\blkda[3]{$(-\infty,0)$}. 
    \begin{jieda}
    令$t=2^x\in(0,+\infty)$, 故方程可化为$t^2+(m+1)t+m=0$. \\
    原问题等价于$t^2+(m+1)t+m=0$在$(0,+\infty)$上只有一根. \\
    而$(t+1)(t+m)=0$, 故$m\in(-\infty,0)$.  
    \end{jieda}
    \item 设非零实常数$a,b,c$满足$a,b$同号, $b,c$异号, 则关于$x$的方程$a\cdot4^x+b\cdot2^x+c=0$\dotfill(\kh{B})
    \xx{无实根}
    {仅有一个实根}
    {有两个异号的实根}
    {有两个同号的实根}
    \begin{jieda}
    $ac<0$, $2^x>0$. 
    \end{jieda}

    \item 若$\alpha,\beta$是方程$\lg^2x-\lg{x^2}-2=0$的两个根, 那么$\log_{\alpha}\beta+\log_{\beta}\alpha$的值为\blkda{$-4$}. 
    
    \begin{jieda}
    令$\lg\alpha=x_1$, $\lg\beta=x_2$, 则$\log_{\alpha}\beta+\log_{\beta}\alpha=\dfrac{\lg\beta}{\lg\alpha}+\dfrac{\lg\alpha}{\lg\beta}=\dfrac{x_2}{x_1}+\dfrac{x_1}{x_2}=\dfrac{(x_1+x_2)^2}{x_1x_2}-2$. \\
    原方程等价于$\lg^2x-2\lg{x}-2=0$, 故$x_1,x_2$是方程$x^2-2x-2=0$的两根. \\
    于是$\log_{\alpha}\beta+\log_{\beta}\alpha=\dfrac{2^2}{-2}-2=-4$. 
    \end{jieda}

    \item
    \begin{enumerate}[(1)]
    \item 方程$\log_2(x+2)=x$的实数根的个数为\blkda{2}; 
    \item 已知$\alpha,\beta$分别是方程$\log_2x+x-3=0$和$2^x+x-3=0$的根, 则$\alpha+\beta$的值为\blkda{3}. 
    \end{enumerate}
    \begin{jieda}
    \begin{enumerate}[(1)]
    \item 作图: $y=\log_2(x+2)$, $y=x$, 故有两个交点. 
    \item 令$f(x)=\log_2x+x$, 则$f(x)$在$(0,+\infty)$上严格单调递增. \\
    于是$f(\alpha)=3$, $f(2^{\beta})=\log_22^{\beta}+2^{\beta}=\beta+2^{\beta}=3$, \\
    故$\alpha=2^\beta$, 故$\alpha+\beta=2^{\beta}+\beta=3$. 
    \end{enumerate}
    \end{jieda}

    \item 函数$f(x)=\log_2(1+3^x+9^x\cdot a)$当$x\leq1$时总有意义, 则实数$a$的取值范围是\blkda[2]{$(-\frac{4}{9},+\infty)$}
    \begin{jieda}
    问题等价于: 对于任意$x\leq 1$, $1+3^x+9^x\cdot a>0$恒成立. \\
    于是$a>\dfrac{-3^x-1}{9^x}=-t^2-t$, 其中$t=\left(\dfrac{1}{3}\right)^x\in\left[\dfrac{1}{3},+\infty\right)$. \\
    而$-t^2-t=-(t+\frac{1}{2})^2+\frac{1}{4}$, 在$[\frac{1}{3},+\infty)$上的最大值为$-\frac{4}{9}$, 故$a>-\frac{4}{9}$. 
    \end{jieda}

    % 1995年全国联赛
    \item 记号$\gauss{x}$表示不大于实数$x$的最大整数, 则方程$\lg^2x-\gauss{\lg{x}}-2=0$的解集为\dotfill(\kh{C})
    \xx{$\left\{x\mid\frac{1}{10}\leq x\leq 100\right\}$}
    {$\left\{\frac{1}{10},10^{\sqrt{2}},10^{-\sqrt{2}},10^{\sqrt{3}}\right\}$}
    {$\left\{\frac{1}{10},10^{\sqrt{3}},100\right\}$}
    {$\left\{\frac{1}{10},10^{\sqrt{2}},10^{-\sqrt{2}},10^{\sqrt{3}},100\right\}$}
    \begin{jieda}
    选择题可以代入求解. \\
    令$t=\lg{x}$, 则$t^2-2=\gauss{t}\in(t-1,t]$, 解得$t\in\left[-1,\dfrac{1-\sqrt{5}}{2}\right)\cup\left(\dfrac{1+\sqrt{5}}{2},2\right]$. \\
    当$t\in\left[-1,\dfrac{1-\sqrt{5}}{2}\right)$时, $\gauss{t}=-1$, 故$t^2-2=-1$, 即$t=-1$(1已舍去); \\
    当$t\in\left(\dfrac{1+\sqrt{5}}{2},2\right)$时, $\gauss{t}=1$, 故$t^2-2=1$, 即$t=\sqrt{3}$($-\sqrt{3}$已舍去); \\
    当$t=2$时, 满足$t^2-2=\gauss{t}$. \\
    综上, $t=-1,\sqrt{3},2$, 故$x=\frac{1}{10},10^{\sqrt{3}},100$. 
    \end{jieda}

    \item 解方程: 
    \begin{enumerate}[(1)] 
    \item $\log_4(13-3x)\cdot\log_{(x-1)}2=1$; 
    \item $\log_2(4^x+4)=x+\log_2(2^{x+1}-3)$. 
    \item $x^{\lg{x}}=\frac{x^3}{100}$; 
    \item $10^{\lg^2x}+x^{\lg{x}}=20$. 
    \end{enumerate}
    \begin{jieda}
    \begin{enumerate}[(1)]
    \item $\iff\log_42\log_{(x-1)}(13-3x)=1\iff\log_{(x-1)}(13-3x)=2$. \\
    故$(x-1)^2=13-3x\land x-1>0\land x-1\neq 1$, 解得$x=3$. 
    
    \item $\iff\log_2(4^x+4)=\log_22^x+\log_2(2^{x+1}-3)\iff 4^x+4=2^x(2^{x+1}-3)$. \\
    于是$(2^x+1)(2^x-4)=0$, 故$2^x=4$, 即$x=2$. 
    \item 两边取常用对数, 得$\lg^2{x}=3\lg{x}-2$, 解得$\lg{x}=1$或$2$, 即$x=10$或100. 
    
    \item 首先$x>0$, 于是$x=10^{\lg{x}}$. \\
    故方程可化为$10^{\lg^2x}+(10^{\lg{x}})^{\lg{x}}=20$, 即$2\cdot10^{\lg^2x}=20$. \\
    故$\lg^2x=1$, 即$x=\frac{1}{10}$或10. 
    \end{enumerate}
    \end{jieda}
    \begin{vspace_}
        \vsp[3]
    \end{vspace_}
    \item 已知函数$f(x)=3^x$, 且$f(a+2)=18$, 函数$g(x)=3^{ax}-4^x$的定义域为$[-1,1]$. 
    \begin{enumerate}[(1)]
    \item 求函数$g(x)$的解析式; 
    \item 若方程$g(x)=m$有解, 求实数$m$的取值范围.
    \end{enumerate}
    \begin{jieda}
    \begin{enumerate}[(1)]
    \item $3^{a+2}=18$, 则$3^a=2$. \\
    于是$g(x)=(3^a)^x-4^x=2^x-4^x$. 
    \item $x\in[-1,1]$时, $2^x\in[\frac{1}{2},2]$, 于是$g(x)\in[-2,\frac{1}{4}]$, 故$m\in[-2,\frac{1}{4}]$. 
    \end{enumerate}
    \end{jieda}
    \begin{vspace_}
        \vsp[3]
    \end{vspace_}

    \item 若方程$\lg(2x)\lg(3x)+a^2=0$有两个不相等的实数根, 求实数$a$的取值范围, 并求方程的两根的积. 
    
    \begin{jieda}
    $\iff -a^2=(\lg2+\lg{x})(\lg3+\lg{x})\iff -a^2=\lg^2x+\lg6\cdot\lg{x}+\lg2\lg3$. \\
    于是$-a^2>-\left(\frac{\lg\frac{3}{2}}{2}\right)^2$, 故$-\frac{\lg\frac{3}{2}}{2}<a<\frac{\lg\frac{3}{2}}{2}$. \\
    此时$\lg{x_1}+\lg{x_2}=-\lg6$, 即$x_1x_2=\frac{1}{6}$. 
    \end{jieda}
    \begin{vspace_}
        \vsp[3]
    \end{vspace_}
    
    \item 关于$x$的方程$\log_2(x+\sqrt{x^2-1})-a=0$有实数解, 求实数$a$的取值范围, 并求出方程的解.
    \begin{jieda}
    $x+\sqrt{x^2-1}>0\iff x\geq 1$, 于是$\log_2(x+\sqrt{x^2-1})\geq 0$, 故$a\geq 0$. \\
    又$2^a=x+\sqrt{x^2-1}\iff 2^a=\dfrac{1}{x-\sqrt{x^2-1}}$, \\
    于是$\begin{cases}
    2^a=x+\sqrt{x^2-1}\\
    2^{-a}=x-\sqrt{x^2-1}
    \end{cases}$, 
    解得$x=\dfrac{2^a+2^{-a}}{2}$. 
    \end{jieda}
    \begin{vspace_}
        \vsp[3]
    \end{vspace_}
    
    \item 已知函数$f(x)=\lg(10^x-1)$, 如果方程$f^{-1}(2x)=\lambda+f(x)$有实根, 求实数$\lambda$的取值范围. 
    \begin{jieda}
    $f^{-1}(x)=\lg(10^x+1)$, 于是$\lambda=f^{-1}(2x)-f(x)=\lg\dfrac{10^{2x}+1}{10^x-1}(x>0)$. \\
    即$10^{\lambda}=\dfrac{10^{2x}+1}{10^x-1}=10^x-1+\dfrac{2}{10^x-1}+2\geq 2\sqrt{2}+2$, 故$\lambda\geq\lg(2\sqrt{2}+2)$. 
    \end{jieda}
    \begin{vspace_}
        \vsp[3]
    \end{vspace_}
    
    \item 关于$x$的方程$k\cdot9^x-k\cdot3^{x+1}+6(k-5)=0$在$[0,2]$上有唯一解, 求实数$k$的取值范围. 
    \begin{jieda}
    令$t=3^x\in[1,9]$, 于是关于$t$的方程$kt^2-3kt+6(k-5)=0$在$[1,9]$上有唯一解. \\
    于是$k=\dfrac{30}{t^2-3t+6}\iff\dfrac{30}{k}=t^2-3t+6$. \\
    结合图像可得$\dfrac{30}{k}\in(4,60]\cup\{\frac{15}{4}\}$, 即$k\in[\frac{1}{2},\frac{15}{2})\cup\{8\}$. 
    \end{jieda}
    \begin{vspace_}
        \vsp[3]
    \end{vspace_}

    \item 如果方程$\log_a\dfrac{x-5}{x+5}=1+\log_a(x-3)$有实根, 求实数$a$的范围. 
    \begin{jieda}
    $\dfrac{x-5}{x+5}>0\land x-3>0$, 则$x\in(5,+\infty)$. \\
    于是方程可化为$\dfrac{x-5}{x+5}=a\cdot(x-3)$, 即$\dfrac{1}{a}=\dfrac{(x+5)(x-3)}{x-5}$. \\
    $\dfrac{1}{a}=x-5+\dfrac{20}{x-5}+12\geq4\sqrt{5}+12$, 且$a>0,a\neq 1$, 故$a\in\left(0,\dfrac{3-\sqrt{5}}{16}\right]$. 
    \end{jieda}
    \begin{vspace_}
        \vsp[3]
    \end{vspace_}
    
    \item 实数$a$取何值时, 方程$\lg(x-1)+\lg(3-x)=\lg(1-ax)$有一解? 两解? 无解? 
    \begin{jieda}
    $x\in(1,3)$, 于是$(x-1)(3-x)=1-ax$, 即$a=x+\dfrac{4}{x}-4$. \\
    结合图像分析易得
    当$a\in\{0\}\cup[\frac{1}{3},1)$时, 只有一解; \\
    当$a\in(0,\frac{1}{3})$时, 有两解; \\
    其他情形, 无解. 
    \end{jieda}
    \begin{vspace_}
        \vsp[3]
    \end{vspace_}

    \item 若关于$x$的方程$\lg(ax)\cdot\lg(ax^2)=4$有两个小于1的正根$\alpha,\beta$, 
    且满足$\abs{\lg\alpha-\lg\beta}\leq 2\sqrt{3}$, 求实数$a$的取值范围. 
    \begin{jieda}
    $ax>0\land ax^2>0$, 故$a>0,x>0$. \\
    于是原方程可化为$(\lg{a}+\lg{x})(\lg{a}+2\lg{x})=4$. \\
    令$t=\lg{x}$, 则$2t^2+3t\lg{a}+\lg^2a-4=0$. \\
    当$x\in(0,1)$时, $t<0$. \\
    于是$\begin{cases}
    \Delta\geq 0\\
    t_1+t_2<0\\
    t_1t_2>0\\
    \abs{t_1-t_2}\leq 2\sqrt{3}
    \end{cases}$, 
    解得$a\in(100,10000]$. 
    \end{jieda}
    \begin{vspace_}
        \vsp[3]
    \end{vspace_}

    \item 已知关于$x$的方程$x^2-5xlog_2a+6(log_2a)^2 = 0$有实根，其中仅一个较小的根在区间$(1, 2)$内，求$a$的取值范围.
    \begin{jieda}
        $x^2-5xlog_2a+6(log_2a)^2 = 0$即$(x - 2log_2a)(x - 3log_2a) = 0$，故$1 < 2log_2a < 2 \leq 3log_2a$
        因此$log_2a \in [\frac{2}{3}, 1)$，解得$a \in [\sqrt[3]{4}, 2)$
    \end{jieda}
    \begin{vspace_}
        \vsp[3]
    \end{vspace_}

    \item 已知关于$x$的方程$3^{2x+1}+(m-1)(3^{x+1}-1)-(m-3)3^x=0$有两个不同的实数根, 求实数$m$的取值范围. 
    \begin{jieda}
    令$t=3^x\in(0,+\infty)$, 于是问题转化为关于$t$的方程$3t^2+2mt+1-m=0$有两个不等正根. \\
    故$\begin{cases}
    \Delta>0\\
    t_1+t_2>0\\
    t_1t_2>0
    \end{cases}$, 解得$m<\dfrac{-3-\sqrt{21}}{2}$. 
    \end{jieda}
    \begin{vspace_}
        \vsp[3]
    \end{vspace_}
\end{enumerate}

\section*{指对方程(2)}
\begin{enumerate}
    \item 

\begin{enumerate}
\item 方程$\log_4(2-x)=\log_2(x-1)-1$的解是\blkda{$2\sqrt{2}-1$}. 
\begin{jieda}
注意到$x\in(1,2)$, 故原方程
$\iff\dfrac{1}{2}\log_2(2-x)=\log_2(x-1)-1$, \\
即$\log_2(2-x)+2=2\log_2(x-1)$. \\
于是$4(2-x)=(x-1)^2\land x\in(1,2)$, 解得$x=2\sqrt{2}-1$. 
\end{jieda}

\item 方程$3^{\sqrt{5x-2y}}+5^{\sqrt{y-5}}=2$的解为\blkda[2]{$x=2,y=5$}. 
\begin{jieda}
$3^{\sqrt{5x-2y}}\geq 1, 5^{\sqrt{y-5}}\geq 1$, 故原方程等价于$3^{\sqrt{5x-2y}}=5^{\sqrt{y-5}}=1$. \\
解得$x=2,y=5$. 
\end{jieda}

\item 方程$\log_a(x+2)=\sqrt{-x}(a>0, a\neq 1)$的实数解的个数为\blkda{1}. 
\begin{jieda}
作图, 当$a\in(0,1)$时, 由凹凸性只有一解; $a>1$时, 只有一解(可由单调性证得). 
\end{jieda}


\item 方程$5^{5x+1}=3^{x-1}$的解集是\blkda[3]{$\left\{\dfrac{\lg15}{\lg3-5\lg5}\right\}$}. 


\item 设$x,y$是实数, 且满足$\begin{cases}
(x-1)^3+1997(x-1)=-1\\
(y-1)^3+1997(y-1)=1
\end{cases}$, 则$x+y$的值是\blkda{2}.
\begin{jieda}
令$f(t)=t^3+1997t$, 则$f(t)$是$\bR$上的奇函数, 且为严格增函数. \\
又$f(x-1)=-1\land f(y-1)=1$, 故$x-1+y-1=0$, 即$x+y=2$. 
\end{jieda}




\item 若关于$x$的不等式$\abs{x-1}>x^2+a$仅有负数解, 则实数$a$的取值范围是\blkda{$(1,\frac{5}{4})$}. 
\begin{jieda}
$\iff \abs{x-1}-x^2>a$. 作$y=\abs{x-1}-x^2$的草图如下: \\
\begin{minipage}{\linewidth}
\centering
\includegraphics[totalheight=3.5cm]{./bodyxb/src/main42chap21zhiduishufangcheng02.png}
\end{minipage}
结合图像可得$a\in(1,\frac{5}{4})$. 
\end{jieda}



\item $f(x)=x^3+bx+c$是$[-1,1]$上的增函数, 且$f(-\frac{1}{2})\cdot f(\frac{1}{2})<0$, 
则$f(x)=0$在$[-1,1]$内\dotfill(\kh{C})
\xx{可能有3个实根}{可能有2个实根}{有唯一实根}{没实根}
\begin{jieda}
由条件可知, 在$[-1,1]$上严格单调递增. \\
故当$-1\leq x_1<x_2\leq 1$时, $\dfrac{f(x_2)-f(x_1)}{x_2-x_1}=x_2^2+x_2x_1+x_1^2+b>0$恒成立, 
即$b>-x_1^2-x_1x_2-x_2^2$恒成立. \\
而$x_1^2+x_1x_2+x_2^2\in(0,3)$, 故$b\geq 0$. \\
易得$f(x)$在$\bR$上严格单调递增. (和函数的单调性)
\end{jieda}


\item 方程$\lg{x}+x-3=0$的根所在的区间是\dotfill(\kh{B})
\xx{$({1,2})$}{$(\frac{5}{2},\frac{11}{4})$}{$(\frac{9}{4},\frac{5}{2})$}{$(3,\frac{13}{4})$}
\begin{jieda}
$f(x)=\lg{x}+x-3$在$(0,+\infty)$上严格单调递增, 且$f(\frac{5}{2})f(\frac{11}{4})<0$. 
\end{jieda}


\item 设函数$y=f(x)$的反函数为$y=f^{-1}(x)$, 且$y=f(2x-1)$的图象过点$(\frac{1}{2},1)$, 则$y=f^{-1}(x)$的图象必过点\dotfill(\kh{C})
\xx{$(\frac{1}{2},1)$}{$(1,\frac{1}{2})$}{$(1,0)$}{$(0,1)$}
\begin{jieda}
$f(2\cdot\frac{1}{2}-1)=1\iff f(0)=1\iff 0=f^{-1}(1)$. 
\end{jieda}

% [source={1999年全国联赛}]
\item 已知$(\log_23)^x-(\log_53)^x\geq(\log_23)^{-y}-(\log_53)^{-y}$, 则\dotfill(\kh{B})
\xx{$x-y\geq 0$}{$x+y\geq 0$}{$x-y\leq 0$}{$x+y\leq 0$}
\begin{jieda}
令$f(t)=(\log_23)^t-(\log_53)^t$, 则$f(t)$在$\bR$上严格单调递增. \\
于是$f(x)\geq f(-y)\iff x\geq -y$, 即$x+y\geq 0$. 
\end{jieda}


\item 已知关于$x$的方程$\log_a(2x^2+x-3)-\log_a(x+4)=1+\log_a(a-1)$有实数解, 且其中之一在2与3之间. 
\begin{enumerate}[(1)]
\item 求实数$a$的取值范围; 
\item 若$a$是整数, 求方程的解.
\end{enumerate}
\begin{jieda}
$2x^2+x-3>0\land x+4>0\iff x\in(-4,-\frac{3}{2})\cup(1,+\infty)$. \\
$a-1>0\land a\in(0,1)\cup(1,+\infty)\iff a>1$. 
\begin{enumerate}[(1)]
\item $\iff \log_a\dfrac{2x^2+x-3}{x+4}=\log_aa(a-1)$. \\
故$a(a-1)=\dfrac{2x^2+x-3}{x+4}=2(x+4)+\dfrac{25}{x+4}-15$. \\
当$x\in(2,3)$时, $x+4\in(6,7)$, 故$2(x+4)+\dfrac{25}{x+4}-15\in(\frac{7}{6},\frac{18}{7})$. \\
于是$a(a-1)\in(\frac{7}{6},\frac{18}{7})\land a>1$, 故$a\in\left(\dfrac{3+\sqrt{51}}{6},\dfrac{7+\sqrt{553}}{14}\right)$. 

\item $a(a-1)\in(\frac{7}{6},\frac{18}{7})\land a>1\land a\zz$, 故$a=2$. \\
于是$2=\dfrac{2x^2+x-3}{x+4}$, 解得$x=\dfrac{1\pm\sqrt{89}}{4}$. 检验后发现符合条件. 
\end{enumerate}
\end{jieda}


\item 已知关于$x$的方程$(2^x-1)^2-2\cdot\abs{2^x-1}+k=0$有三个根, 求实数$k$的取值范围. 
\begin{jieda}
令$t=\abs{2^x-1}$, 该方程的根的情形如下: \\
$t<0$时, 无根; $t\in\{0\}\cup[1,+\infty)$时, 一根; $t\in(0,1)$时, 两根. \\
于是对于$t^2-2t+k=0$而言, $t_1\in\{0\}\cup[1,+\infty), t_2\in(0,1)$. \\
又$t_1+t_2=2$, 故$t_1\in(1,+\infty),t_2\in(0,1)$. \\
令$f(t)=t^2-2t+k$, 于是只要$f(0)>0\land f(1)<0$即可, 解得$k\in(0,1)$. 
\end{jieda}


\item 解关于$x$的方程$\lg{x}+\lg(4-x)=\lg(a+2x)$, 并讨论方程解的个数. 
\begin{jieda}
注意到$x\in(0,4)$, 于是$x(4-x)=a+2x$, 即$a=2x-x^2$. 
作$y=2x-x^2, x\in(0,4)$的草图, 分析得如下结果: \\
$a\in(-\infty,-8]\cup(1,+\infty)$时, 无解; \\
$a\in(-8,0]\cup\{1\}$时, 一解; \\
$a\in(0,1)$时, 两解. 
\end{jieda}


\item 
\begin{enumerate}[(1)]
\item 阅读不等式$2^x+1>3^x$的解法: \\
设$f(x)=(\frac{2}{3})^x+(\frac{1}{3})^x$, 则$f(x)$在$\bR$上单调递减. \\
因为$f(1)=1$, 所以当$x<1$时, $(\frac{2}{3})^x+(\frac{1}{3})^x>1$; 当$x\geq 1$时, $(\frac{2}{3})^x+(\frac{1}{3})^x\leq 1$.\\
因为$3^x>0$, 所以不等式$2^x+1>3^x$的解集为$\left\{x\mid x<1\right\}$. \\
试利用上面的解法解不等式$2^x+3^x\geq 5^x$; 

\item 证明: $5^x+12^x=13^x$有且只有一个实数解$x=2$. 
\end{enumerate}
\begin{jieda}
\begin{enumerate}[(1)]
\item $\iff(\frac{2}{5})^x+(\frac{3}{5})^x\geq 1$. \\
令$f(x)=(\frac{2}{5})^x+(\frac{3}{5})^x$, 则$f(x)$在$\bR$上严格单调递减, 且$f(1)=1$, \\
故$f(x)\geq 1$的解为$(-\infty,1]$. 

\item 令$f(x)=(\frac{5}{13})^x+(\frac{12}{13})^x$, 则$f(x)$在$\bR$上严格单调递减, 且$f(2)=1$. \\
原方程等价于$f(x)=1$, 故方程只有一解$x=2$. 
\end{enumerate}
\end{jieda}


\end{enumerate}

\end{enumerate}


 \section*{对数函数与反函数阶段练习}

\begin{enumerate}
\item 函数$y=2^{1-x}+3$的反函数的解析式为\blkda[4]{$y=1-\log_2(x-3)$}. 

\item 若$\log_m\frac{2}{5}<1$, 则实数$m$的取值范围是\blkda[3]{$(0,\frac{2}{5})\cup(1,+\infty)$}. 

\item 设函数$f(x)=\begin{cases}
2^{-x}, & x\leq 1\\
\log_{81}x, & x>1
\end{cases}$, 
则满足$f(x)=\frac{1}{4}$的$x$的值为\blkda{$3$}. 
\begin{jieda}
条件方程, 分类即可. 
\end{jieda}


\item 函数$y=\sqrt{\log_{0.5}(4x^2-3x)}$的定义域为\blkda[3]{$[-\frac{1}{4},0)\cup(\frac{3}{4},1]$}. 
\begin{jieda}
$\log_{0.5}(4x^2-3x)\geq0\iff 0<4x^2-3x\leq 1$. 
\end{jieda}


\item 函数$f(x)=\log_2\left(x-\frac{a}{x}\right)(a>0)$在区间$[2,+\infty)$上是增函数, 则$a$的取值范围是\blkda{$(0,4)$}. 
\begin{jieda}
当$a>0$时, $y=x-\frac{a}{x}$在$(0,+\infty)$上单调递增. \\
考虑到定义域, 于是$2-\frac{a}{2}>0$, 故$a\in(0,4)$. 
\end{jieda}


\item 函数$y=\abs{x^2-1}+1$的图像与函数$y=\lg(x+10)$的图像交点个数为\blkda{2}. 



\item 已知函数$y=\log_{\frac{1}{2}}(x^2-ax-a)$在区间$(-\infty,1-\sqrt{3})$上递增, 则$a$的取值范围是\blkda{$[2-2\sqrt{3},2]$}. 
\begin{jieda}
$\frac{a}{2}\geq 1-\sqrt{3}\land (1-\sqrt{3})^2-a(1-\sqrt{3})-a\geq 0$. 
\end{jieda}

\item 设$f(x)$是奇函数, 当$x>0$时, $f(x)=\log_2{x}$, 则当$x<0$时, $f(x)$的解析式为\blkda[3]{$f(x)=-\log_2(-x)$}. 

\item 当$x\in(1,2)$时, 不等式$(x-1)^2<\log_ax$恒成立, 则$a$的取值范围是\blkda{$(1,2]$}. 
\begin{jieda}
数形结合, 讨论: $a>1\land \log_a2\geq (2-1)^2$. 
\end{jieda}


\item 若函数$f(x)=\left(\log_2\dfrac{x}{2}\right)\left(\log_2\dfrac{x}{4}\right)$的定义域是不等式$2(\log_{\frac{1}{2}}x)^2+7(\log_{\frac{1}{2}}x)+3\leq 0$的解集, 
则函数$f(x)$的最大值为\blkda{2}, 最小值为\blkda{$-\frac{1}{4}$}. 
\begin{jieda}
参见\autoref{subsec:secduishuhanshudexingzhi}\autoref{item:secduishuhanshudexingzhi13}
\end{jieda}

\item 如果$\log_a2>\log_b2$, 那么$a,b$间的关系是\dotfill(\kh{D}) 
\xx{$b>a>1$}
{$1>b>a>0$}
{$a>1$且$0<b<1$}
{以上三种答案均有可能}
\begin{jieda}
$\iff\dfrac{\lg2}{\lg{a}}>\dfrac{\lg2}{\lg{b}}\iff\dfrac{1}{\lg{a}}>\dfrac{1}{\lg{b}}$. 
\end{jieda}


\item  函数$y=\abs{\lg\abs{x}}$的图像是\dotfill(\kh{B})
\begin{minipage}{\linewidth}
\centering
\includegraphics[totalheight=3cm]{./bodyxb/src/dasfnashgkjhgkhrg.jpg}
\end{minipage}

\begin{jieda}
\begin{minipage}{\linewidth}
\centering
\includegraphics[totalheight=3cm]{./bodyxb/src/main42chap21duishuhanshuzonghe01.png}
\end{minipage}
\end{jieda}

\item 与函数$y=x$为同一函数的是\dotfill(\kh{D})
\xx{$y=x\log_ax$}
{$y=\sqrt{x^2}$}
{$y=a^{\log_ax}(a>0, a\neq1)$}
{$y=\log_aa^x({a>0, a\neq1})$}
\begin{jieda}
定义域. 
\end{jieda}


\item 已知函数$f(x)=\log_a(x-k)$的图像过点$(4,0)$而且其反函数$f^{-1}(x)$的图像过点$(1,7)$, 则$f(x)$是\dotfill(\kh{A})
\xx{增函数}{减函数}{奇函数}{偶函数}
\begin{jieda}
$k=3,a=4$. 
\end{jieda}


\item 右图是函数$y=a+\log_bx$的图像, 其中$a,b$为常数, 则下列结论正确的是\dotfill(\kh{D})
\xx{$a<0,b>1$}{$a>0,b>1$}{$a>0,0<b<1$}{$a<0,0<b<1$}
\begin{minipage}{\linewidth}
\centering
\includegraphics[totalheight=3cm]{./bodyxb/src/ejrgfjnjanv.jpg}
\end{minipage}





\item 对于函数$y=\lg\frac{x}{100}$的图像给出三个命题: 
\begin{dingautolist}{172}
\item 存在直线$l_1$, 函数$y=\lg\frac{x}{100}$的图像与函数$y=100\cdot10^x$的图像关于直线$l_1$对称;
\item 存在直线$l_2$, 函数$y=\lg\frac{x}{100}$的图像与函数$y=\log_{0.1}\frac{x}{100}$的图像关于直线$l_2$对称; 
\item 存在直线$l_3$, 函数$y=\lg\frac{x}{100}$的图像与函数$y=\lg(-x)+\log_{0.1}100$的图像关于直线$l_3 $对称. 
\end{dingautolist}
上述命题中正确命题的个数是\dotfill(\kh{A})
\xx{3}{2}{1}{0}
\begin{jieda}
令$f(x)=\lg\frac{x}{100}$. 
\begin{dingautolist}{172}
\item 后者是前者的反函数$y=f^{-1}(x)$. 
\item 后者是$y=-f(x)$. 
\item 后者是$y=f(-x)$. 
\end{dingautolist}
\end{jieda}


\item 设函数$f(x)=\log_a\left(1-\dfrac{a}{x}\right)$, 其中$0<a<1$. 
\begin{enumerate}[(1)]
\item 证明: $f(x)$是$(a,+\infty)$上的减函数; 
\item 解不等式$f(x)>1$.
\end{enumerate}
\begin{jieda}
\begin{enumerate}[(1)]
\item 内层单调递减. 需要作差的过程. 
\item 结合单调性可得$f(x)>1\iff 0<1-\frac{a}{x}<a\iff 1-a<\frac{a}{x}<1\iff a<x<\frac{a}{1-a}$. 
\end{enumerate}
\end{jieda}



\item 已知函数$f(x)=\lg(1+2^x+3^x+4^x+a\cdot5^x)$对一切$x\leq 2$有意义, 求实数$a$的取值范围. 
\begin{jieda}
$1+2^x+3^x+4^x+a\cdot5^x>0$恒成立, 即$a>-\dfrac{1+2^x+3^x+4^x}{5^x}$恒成立. \\
而$y=-\dfrac{1+2^x+3^x+4^x}{5^x}$在$\bR$上单调递增, 故$a>-\dfrac{6}{5}$($x=2$时). 
\end{jieda}



\item 已知函数$f(x)$满足$f(\log_ax)=\dfrac{a}{a^2-1}\left(x-\dfrac{1}{x}\right)(a>0, a\neq 1)$. 
\begin{enumerate}[(1)]
\item 求$f(x)$的解析式, 并判定奇偶性及单调性; 
\item 当$x\in(-\infty,2)$时, $f(x)-4$的值恰好为负数, 求$a$的取值范围. 
\end{enumerate}
\begin{jieda}
\begin{enumerate}
\item 令$t=\log_ax$, 于是$f(t)=\dfrac{a}{a^2-1}(a^t-a^{-t})$. \\
易得$f(x)$为奇函数, 在$\bR$上严格单调递增. 
\item 结合单调性, 只要$f(2)\leq 4$即可. \\
解得$a\in[2-\sqrt{3},1\cup(1,2+\sqrt{3}]$. 
\end{enumerate}
\end{jieda}


\item 设函数$f(x)=\log_{\frac{1}{3}}\left(kx^2+(k+2)x+k+2\right)(k\zr)$. 
\begin{enumerate}[(1)]
\item 若$f(x)$定义域为$\bR$, 求$k$的取值范围;
\item 若$f(x)$的值域为$\bR$, 求$k$的取值范围;
\item 若$f(x)$有最小值, 求$k$的取值范围.
\end{enumerate}
\begin{jieda}
\begin{enumerate}[(1)]
\item $k=0$不符合条件. \\
$k\neq 0$时, $kx^2+(k+2)x+k+2>0$恒成立, 于是$\Delta<0\land k>0$. 解得$k\in(\frac{2}{3},+\infty)$. 
\item $y=kx^2+(k+2)x+k+2$能遍历$(0,+\infty)$. \\
$k=0$可以; $k\neq 0$时, $\Delta\geq 0\land k>0$, 故$k\in(0,\frac{2}{3}]$. \\
综上, $k\in[0,\frac{2}{3}]$. 
\item $y=kx^2+(k+2)x+k+2$有最大值, 且最大值为正. \\
$k=0$不符合条件; $k\neq 0$时, $k<0\land \Delta>0$, 解得$k\in(-2,0)$. \\
综上, $k\in(-2,0)$. 
\end{enumerate}
\end{jieda}


\item 已知函数$f(x)=2^x-1$有反函数为$f^{-1}(x)$, $g(x)=\log_4(3x+1)$. 
\begin{enumerate}[(1)]
\item 若$f^{-1}(x)\leq g(x)$, 求$x$的取值范围$D$; 
\item 设函数$H(x)=g(x)-\frac{1}{2}f^{-1}(x)$, 当$x\in D$时, 求函数$H(x)$的值域. 
\end{enumerate}
\begin{jieda}
$f^{-1}(x)=\log_2(x+1)$. 
\begin{enumerate}[(1)]
\item $\log_2(x+1)<\log_4(3x+1)\iff 0<x+1<\sqrt{3x+1}$, 故$D=[0,1]$. 

\item $H(x)=\log_4(3x+1)-\frac{1}{2}\log_2(x+1)=\frac{1}{2}(\log_2(3x+1)-\log_2(x+1))=\frac{1}{2}\log_2\frac{3x+1}{x+1}$. \\
易证$H(x)$在$D$上单调递增, 故$H(x)\in[0,\frac{1}{2}]$. 
\end{enumerate}
\end{jieda}


\end{enumerate}